\documentclass[11pt,openany]{book}

\usepackage[paperwidth=6in,paperheight=9in,inner=0.82in,outer=0.68in,top=0.75in,bottom=0.82in]{geometry}
\usepackage{microtype}
\usepackage{setspace}
\usepackage{amsmath,amssymb,amsthm,mathtools}
\usepackage{booktabs}
\usepackage{enumitem}
\usepackage{xcolor}
\usepackage{hyperref}

\definecolor{inkblue}{HTML}{17324D}
\hypersetup{
  colorlinks=true,
  linkcolor=inkblue,
  citecolor=inkblue,
  urlcolor=inkblue,
  pdftitle={The Witness Plane},
  pdfauthor={Flyxion}
}

\setstretch{1.16}
\setlength{\parindent}{1.1em}
\setlength{\parskip}{0.12em}
\setlength{\emergencystretch}{2em}
\pagestyle{plain}

\newtheorem{definition}{Definition}[chapter]
\newtheorem{principle}[definition]{Principle}
\newtheorem{invariant}[definition]{Invariant}
\newtheorem{proposition}[definition]{Proposition}
\newtheorem{lemma}[definition]{Lemma}
\newtheorem{theorem}[definition]{Theorem}
\newtheorem{assumption}[definition]{Assumption}
\newtheorem{failuremode}[definition]{Failure Mode}

\newcommand{\Pop}{\mathsf{POP}}
\newcommand{\Refuse}{\mathsf{REFUSE}}
\newcommand{\Bind}{\mathsf{BIND}}
\newcommand{\Transform}{\mathsf{TRANSFORM}}
\newcommand{\Verify}{\mathsf{VERIFY}}
\newcommand{\Collapse}{\mathsf{COLLAPSE}}
\newcommand{\Ledger}{\mathsf{LEDGER}}
\newcommand{\Hash}{\mathsf{Hash}}
\newcommand{\Pass}{\mathsf{PASS}}
\newcommand{\Fail}{\mathsf{FAIL}}
\newcommand{\Indeterminate}{\mathsf{INDETERMINATE}}
\newcommand{\ErrorVerdict}{\mathsf{ERROR}}
\newcommand{\Waived}{\mathsf{WAIVED}}
\newcommand{\Required}{\mathsf{REQUIRED}}
\newcommand{\interlude}{\cleardoublepage}

\title{\textbf{The Witness Plane}\\[5pt]
\large Admissibility, Verification, and Commitment in Event Systems}
\author{Flyxion\\\small Independent Researcher}
\date{September 2026}

\begin{document}

\frontmatter
\maketitle

\chapter*{Orientation}
\addcontentsline{toc}{chapter}{Orientation}

An event system does not become trustworthy merely because it preserves an
append-only log, executes deterministic code, or attaches hashes to its
outputs. These mechanisms establish that something was recorded or computed;
they do not establish that the input was admissible, that the computation was
authorized, that its obligations were satisfied, or that its result was fit to
alter the operational world. This book presents Spherepop as an
architecture for keeping those propositions separate.

The central pipeline is
\[
\begin{aligned}
  \Pop &\longrightarrow \Refuse \longrightarrow \Bind \\
       &\longrightarrow \Transform \longrightarrow \Verify
        \longrightarrow \Collapse,
\end{aligned}
\]
while \(\Ledger\) is reinterpreted as the persistent witness plane beneath all
six transitions. The ledger records encounters, refusals, bindings,
executions, verdicts, failed commitments, successful commitments, and recovery
operations. Operational history advances only when a verified candidate
satisfies the policy-relative collapse guard.

Its philosophical thesis is that occurrence, memory, admissibility,
execution, verification, and commitment are different achievements. Its
technical thesis is that a system becomes more inspectable when each
achievement is represented by a typed transition, a declared policy, a
versioned witness, and a distinct failure vocabulary. Its practical thesis is
that distributed systems should target accountable commitment rather than the
fiction of absolute warrant.

\section*{Composition and method}

The book comprises thirty-five chapters, four interludes, and eight technical
appendices. Its argument moves from conceptual separations through a typed
transition system, then into concurrency, external effects, recovery,
implementation, evaluation, and institutional consequence. Formal chapters
pair definitions and positive results with counterexamples, implementation
interpretations, and explicit statements of what the machinery does not
establish. The appendices consolidate the notation, operational semantics,
proof obligations, serialization rules, example records, recovery cases,
security boundaries, and open research programme needed to make those claims
inspectable as a whole.

\section*{Canonical status}

The Spherepop pipeline and its event semantics are the principal subject of
the book and form one focused elaboration of the author's broader
Spherepop research programme \cite{flyxion-spherepop}. Connections to RSVP,
Kinetic-Event Synthesis, Distinction
Holonomy Theory, and MEM|8 are presented as bounded integration hypotheses,
not as foundations on which the core architecture depends. Detailed KES
mathematics not independently established in its source manuscripts remains
explicitly provisional. The architecture specified here is substantially
complete within its declared scope, while the extensions collected in the
research programme remain available for later development without being
presupposed by the results already stated.

\tableofcontents

\mainmatter

\part{The Separations}

\chapter{The Event Is Not Yet an Action}
\section{Chapter thesis}
Most systems that process events collapse a long sequence of distinct
achievements into a single word: ``handling'' the event. A request arrives, is
parsed, is executed, is checked, and its effect becomes part of the system's
permanent record---and all of this is described, and often implemented, as one
operation. The request ``was handled.'' The claim ``was processed.'' The
instruction ``went through.''

This chapter argues that the compression is a category error, and that the
error is not merely stylistic. Receiving an event, recording that it occurred,
parsing its structure, executing a computation over it, validating the result
against a standard, and committing that result to the system's operational
history are six different achievements. Each can succeed while the others
fail. A malformed request can be received and recorded without being parsed. A
request can be parsed and executed without its output ever being checked. An
output can be checked and found wanting without anything having been
committed. And, critically, an output can be checked, found acceptable, and
still not be the thing that was committed, if the committing step observes a
world that changed between verification and commitment.

Treating these as one operation does not make a system simpler. It makes a
system's failures unintelligible, because there is no vocabulary left for
describing which of the six achievements actually happened when something
goes wrong. This book exists to supply that vocabulary, and this chapter
exists to motivate why it is needed before any formal machinery is introduced.

The six achievements are an introductory decomposition rather than a premature
one-to-one transcription of the Spherepop operators. Later chapters refine the
space between encounter and execution through REFUSE and BIND, and show that
LEDGER is not merely another sequential achievement but the witness plane
beneath them all.

\section{Core distinctions}

\subsection{Occurrence versus endorsement}

An event can occur---arrive at a system's boundary, be observed, be
timestamped---without the system endorsing it in any sense stronger than
``this was seen.'' Occurrence is a fact about the world: something crossed a
boundary at a given time, from a given source, with a given content.
Endorsement is a fact about the system: it has decided the occurrence is fit to
influence what happens next. Confusing the two produces the naive assumption
that anything a system has recorded is something the system has approved. It
is not. A system can, and generally should, keep a durable record of things it
explicitly declines to approve.

\subsection{Evidence versus warrant}

A closely related distinction concerns what an occurrence is good for. An
occurrence can serve as evidence---as something that bears on a later
decision---long before it constitutes a warrant for that decision. A single
unconfirmed report is evidence that something may be true. It is not, by
itself, a warrant for acting as though it is true. Systems that conflate
evidence with warrant tend to act on the first plausible signal they receive;
systems that keep the distinction can accumulate evidence, let it mature, and
act only once a declared threshold of warrant has been met.

\subsection{Computation versus authorization}

A computation can execute correctly---produce exactly the output its code
specifies, given its input---without having been authorized to run at all, or
without its output being authorized to take effect. This distinction matters
most where it is most often ignored: in the assumption that if a piece of code
ran successfully, whatever it produced must be legitimate. Correct execution
answers the question ``did the code do what it says.'' It does not answer the
question ``was the code allowed to do this, to this input, on this system's
behalf.'' Those are independent questions, and a system that only asks the
first one has no way to refuse a well-executed but unauthorized computation.

\subsection{History versus audit}

Finally, this book distinguishes two kinds of record-keeping that are
frequently merged into one log. \textbf{Operational history}, written
\(H_t\), is the record of what the system has committed to---the sequence of
transitions that define its current authoritative state. \textbf{Audit
history}, written \(L_t\), is the broader record of everything the system has
encountered, attempted, refused, verified, failed to verify, and eventually
committed or declined to commit. Every element of \(H_t\) should correspond to
something in \(L_t\); the converse should not hold. Most of what a
well-instrumented system remembers is not, and should not be, part of what it
has committed to.

These two histories are formalized fully beginning in Chapter 4. For now, it
is enough to notice that a system with only one log---one that does not
distinguish ``this happened'' from ``this is now authoritative''---cannot
represent the difference between an attempt and a success without inventing ad
hoc conventions for doing so after the fact.

\section{Principal examples}

The four distinctions above are abstract until they are seen doing work. The
following four cases, drawn from different domains, show what goes wrong when
occurrence, evidence, correct execution, and historical record are each
mistaken for something stronger than they are.

\subsection{A malformed network request}

A server receives a request with an invalid content-length header. The request
occurred---it arrived, it consumed a connection, it can be logged. It was never
parsed, because parsing requires a well-formed envelope. A system that logs the
request as ``received'' and stops there has correctly distinguished occurrence
from endorsement. A system that treats ``received'' as equivalent to ``in the
queue for processing'' will eventually try to process something that cannot be
processed, and will need a second, informal layer of error handling to recover
from a distinction it should have made at intake.

\subsection{A scientifically unsettled claim}

A single study reports an effect. The report is evidence: it should be
retained, cited, and available to later inquiry. It is not a warrant for
treating the effect as established, prescribing based on it, or retracting
prior guidance. The distinction between evidence and warrant is exactly what
allows a claim's status to change over time---as further studies arrive---without
requiring anyone to have acted prematurely on the first report, and without
requiring the first report to be deleted once it is superseded. Chapter 2
returns to this case in detail as the paradigm instance of being admissible to
memory before being admissible to action.

\subsection{A stale financial instruction}

An instruction to transfer funds is authorized at the moment it is issued,
against the account balance and policy in effect at that moment. If the
instruction is executed correctly but only after a delay, during which the
balance or policy has changed, correct execution of the original instruction
is not the same as authorization to execute it now. A system that checks
correctness of execution but not freshness of authorization will faithfully
carry out instructions that are no longer warranted, precisely because it has
conflated two questions that must be asked separately and in order.

\subsection{An irreversible actuator command}

A command to fire a valve, release a mechanism, or otherwise change the
physical world cannot be undone by deleting a log entry. Here the
history-versus-audit distinction is sharpest: a system can and should preserve
full audit history of every candidate command it considered, refused, revised,
and eventually issued---but only one of those records corresponds to an actual
physical consequence, and the system's own bookkeeping must make it possible
to tell, without ambiguity, which one that was. Chapter 17 develops the
further complication that even a correctly authorized commit does not
guarantee the physical action it describes was completed exactly once.

Across all four cases, the failure mode is the same: something true at one
stage---occurrence, evidence, correct execution, or bare record---is treated as
though it settled a later, stronger question it was never equipped to settle.

\section{Proof obligations and status}
\textbf{Foundational but non-theorem-bearing.} This chapter's role is to
stabilize vocabulary and motivate the separations that later chapters
formalize. It supplies typed examples and counterexamples---the four cases
above are chosen so that each isolates exactly one collapsed distinction---but
it does not, and should not, attempt to state or prove a formal result. No
transition system has yet been defined; there is nothing yet for a proposition
to be about. The temptation to open a book on formal architecture with an
early theorem, for the sake of appearing rigorous from page one, should be
resisted here specifically: a theorem proved before its terms are defined is
not rigor, it is the appearance of rigor purchased at the cost of actually
defining anything. The formal apparatus begins in Chapter 5, once the event
envelope itself has a definition to be precise about.

\chapter{Admissible to Memory, Forbidden to Action}
\section{Chapter thesis}
A claim can deserve preservation before it deserves any influence over what a
system does. This is the plainest way to state the Reversibility Gate principle
that grounds the present chapter: epistemic admission, letting something into
the record, and operational admission, letting something change what the
system treats as settled, are separate decisions, and a system is entitled to
make one without making the other. The scientific claim of Chapter 1
illustrated this once. This chapter generalizes it and gives it a name.

\section{Three times}

Confusion between memory and action is usually a confusion about time. Three
moments need to stay distinct.

\textbf{Observation time} is when a system first encounters a claim, request,
or signal. \textbf{Decision time} is when the system evaluates that claim
against a current policy and history, and settles what disposition to give it.
\textbf{Action time} is when, if ever, the claim's consequence is allowed to
change something the system or the world depends on.

These three moments can be far apart, and the distance between them is not a
defect to be engineered away. A claim observed today may not reach decision
time until more corroborating evidence exists. A claim decided favorably at
decision time may not reach action time until a further authorization, a
scheduled window, or a resource becomes available. The error this chapter is
built to prevent is treating the elapsed time between observation and decision,
by itself, as sufficient grounds for advancing to action. Maturity of evidence
is necessary for many kinds of action; it is not sufficient for any of them,
because sufficiency also depends on authorization, capability, and the state
of the world at action time, none of which observation time or the passage of
time can supply on its own.

\section{Layered admission}

A single Boolean property, admitted or not admitted, cannot represent what a
well-run system actually does with a claim. This chapter proposes four separate
predicates instead.

\textbf{Epistemic admission}, written \(\mathrm{Ep}\), holds when a claim is
retained in the system's durable record as something worth knowing about,
independent of whether anything further is done with it.

\textbf{Computational admission}, written \(\mathrm{Comp}\), holds when a
claim is cleared to have some computation, deliberation, or transformation
process run over it.

\textbf{Operational admission}, written \(\mathrm{Op}\), holds when a claim,
or the result of a computation over it, is admitted to become part of the
system's authoritative operational state, the state the system treats as
current and relies on for further decisions.

\textbf{Irreversible admission}, written \(\mathrm{Ir}\), holds when a
claim's consequence is admitted to be enacted in a way that has effects the
system's own record cannot later undo, typically because those effects occur
outside the system's record entirely.

A system that only tracks one of these, usually the last one it happens to
notice, cannot express most of the states a real process passes through. The
remainder of this chapter establishes that the four predicates are logically
independent of one another: none is implied by any other, in either direction,
without an additional architectural guarantee doing the work.

\section{Proof obligations and status}
\textbf{Load-bearing counterexample result.} This section proves pairwise
non-implication among \(\mathrm{Ep}\), \(\mathrm{Comp}\),
\(\mathrm{Op}\), and \(\mathrm{Ir}\). For each of the six unordered pairs,
two witnesses are given: one state satisfying the first predicate and not the
second, and one state satisfying the second and not the first. Each witness
states the assumptions under which it is admissible, since these examples are
drawn from general domains rather than from a single governed system, and a
different set of institutional or architectural assumptions can change which
witnesses apply.

\subsection{Pair \texorpdfstring{\((\mathrm{Ep},\mathrm{Comp})\)}{(Ep, Comp)}}

\(\mathrm{Ep}\land\neg\mathrm{Comp}\): a claim is retained purely as
evidence, with no computation ever cleared to run over it. This is the ordinary
case of an evidence-only refusal disposition, assuming the governing policy
treats retention and computational clearance as separately decidable.

\(\mathrm{Comp}\land\neg\mathrm{Ep}\): a computation runs over an input
that is never itself durably retained. A streaming system that computes an
aggregate statistic from raw sensor readings and discards the raw readings
once the statistic is produced admits the input to computation without
admitting it to the durable record, assuming the discard is a deliberate
storage or privacy policy rather than an unintended loss.

\subsection{Pair \texorpdfstring{\((\mathrm{Ep},\mathrm{Op})\)}{(Ep, Op)}}

\(\mathrm{Ep}\land\neg\mathrm{Op}\): the scientific claim from Chapter 1,
retained in the literature but never adopted into settled guidance, assuming
the domain in question distinguishes publication from adoption.

\(\mathrm{Op}\land\neg\mathrm{Ep}\): a default judgment rendered because a
party failed to appear changes a court's operative record without any evidence
about the merits having been admitted, assuming the governing procedure
permits judgment on absence rather than requiring an evidentiary record in
every case.

\subsection{Pair \texorpdfstring{\((\mathrm{Ep},\mathrm{Ir})\)}{(Ep, Ir)}}

\(\mathrm{Ep}\land\neg\mathrm{Ir}\): most retained claims, including the
scientific claim again, never lead to any irreversible action at all.

\(\mathrm{Ir}\land\neg\mathrm{Ep}\): a mechanical overcurrent interlock
trips a circuit breaker directly, producing a physical change the software
record cannot itself undo before any software-level record of the triggering
condition exists, assuming the interlock is intentionally designed to act at
electromechanical speed ahead of, rather than through, any logging path.

\subsection{Pair \texorpdfstring{\((\mathrm{Comp},\mathrm{Op})\)}{(Comp, Op)}}

\(\mathrm{Comp}\land\neg\mathrm{Op}\): a simulation or dry-run mode executes
the full computation over a candidate input and never commits the result, by
design, assuming the system explicitly separates speculative execution from
commitment.

\(\mathrm{Op}\land\neg\mathrm{Comp}\): an authorized administrator directly
edits the operational record to correct a stuck state, without any submitted
claim having been computed over, assuming the governing policy grants a
privileged manual-override path outside the ordinary pipeline. This witness is
exactly the gap that the collapse guard introduced in Chapter 15 is designed
to close inside Spherepop itself; outside a system that enforces such a guard,
nothing prevents it.

\subsection{Pair \texorpdfstring{\((\mathrm{Comp},\mathrm{Ir})\)}{(Comp, Ir)}}

\(\mathrm{Comp}\land\neg\mathrm{Ir}\): an internal report is generated by a
computation and only ever displayed, never triggering any external effect,
assuming display is not itself classified as an irreversible effect in the
domain at hand.

\(\mathrm{Ir}\land\neg\mathrm{Comp}\): a person presses a physical
emergency-stop button on their own judgment, with no computational process
mediating the decision, assuming the domain permits direct human action
outside the computational pipeline, as most safety-critical domains
deliberately do.

\subsection{Pair \texorpdfstring{\((\mathrm{Op},\mathrm{Ir})\)}{(Op, Ir)}}

\(\mathrm{Op}\land\neg\mathrm{Ir}\): an order is marked approved in a
system's authoritative internal state before the corresponding shipment has
actually occurred, corresponding to a completed prepare step awaiting
actuation, assuming the domain separates internal commitment from external
fulfillment.

\(\mathrm{Ir}\land\neg\mathrm{Op}\): an actuator fires and the physical
effect occurs, but the process that would append the corresponding commit to
the operational history crashes before completing, leaving an orphaned effect
the system's own record does not yet reflect, assuming actuation and
record-keeping are not implemented as one atomic step, which Chapter 17 argues
they generally cannot be.

\subsection{Summary and scope}

The twelve witnesses above establish that no predicate among
\(\mathrm{Ep}\), \(\mathrm{Comp}\), \(\mathrm{Op}\), and
\(\mathrm{Ir}\) implies, or is implied by, any other as a matter of
definition alone. This is a claim about the four predicates taken in
isolation, prior to any particular system's governance. It is not a claim that
a well-designed system should tolerate all twelve gaps equally. Several of the
witnesses above, most sharply the manual-override and orphaned-effect cases,
name exactly the kind of gap that later chapters close with an explicit
architectural guard: the collapse guard of Chapter 15 forecloses
\(\mathrm{Op}\land\neg\mathrm{Comp}\) within Spherepop by construction, and
the prepare/actuate/finalize protocol of Chapter 17 gives
\(\mathrm{Ir}\land\neg\mathrm{Op}\) a named, recoverable status rather than
leaving it as silent data loss. The purpose of proving independence here,
before any such guard has been defined, is to make clear that these
implications are achievements of specific architectural choices, not free
consequences of the four predicates existing at all.

\chapter{Against Absolute Warrant}
\section{Chapter thesis}

It is tempting, once a system distinguishes occurrence from endorsement,
evidence from warrant, and epistemic admission from operational admission, to
conclude that whatever finally clears every gate must be correct. This chapter
argues against that conclusion directly. A transition that has been checked
against every obligation the system declared, using a verifier the system
trusts, at a history head the system confirms is current, remains conditional
on the policy, the specification, the evidence, and the trusted computing base
that produced those checks. None of those four can certify itself as an
absolute ground of warrant. A verified transition is warranted. It is not, and
cannot be, warranted absolutely.

This distinction is not a hedge added for intellectual modesty. It is
load-bearing for how the rest of this book treats the word ``verified.'' Later
chapters build an elaborate apparatus for discharging obligations, and it
would be easy to read that apparatus as a machine for producing certainty. It
is not that. It is a machine for producing a specific, inspectable, revisable
kind of sufficiency, and this chapter exists to say precisely what that
sufficiency does and does not include, before the apparatus is built.

\section{Policy-relative sufficiency}

Call a transition's warrant sufficient, relative to a policy version
\(\sigma_t\), a confirmed history head \(h(H_t)\), and a declared verification
regime, when the transition has satisfied every obligation that \(\sigma_t\)
requires, against the state that \(h(H_t)\) identifies, under the specific
procedures the verification regime names. Each of the three qualifiers is
doing real work, and each is a place where the warrant can be sufficient in
exactly the sense defined and still be wrong.

Sufficiency relative to \(\sigma_t\) means sufficiency relative to whatever
obligations someone wrote into that policy. It says nothing about whether the
policy asked the right questions. Sufficiency relative to \(h(H_t)\) means the
check was performed against a specific, agreed state of the world as the
system understood it at that moment. It says nothing about whether the
system's history correctly reflects the world outside it. Sufficiency relative
to a declared verification regime means the named procedures were followed
and produced the stated verdict. It says nothing about whether those
procedures are capable, in principle, of detecting the kind of error that
happens to be present.

Later chapters give each of these qualifiers a full formal treatment:
\(\sigma_t\) and its versioning in the chapters on binding and obligation
generation, \(h(H_t)\) in the chapters on the collapse guard and concurrency,
and the verification regime in the chapters on verifier independence and
correlated failure. This chapter's job is only to fix the shape of the claim
before that machinery exists, so that later precision is read as precision
about a bounded thing, not as a route to an unbounded one.

\section{Residual fallibility}

Five distinct sources of residual fallibility survive even a maximally careful
verification regime. None of them is a flaw in the Spherepop architecture
specifically; each is a property of any system that must decide, using finite
declared rules, whether to trust something.

\subsection{Specification error}

An obligation set can be internally consistent, fully discharged, and still
wrong, because it encoded the wrong requirement. A binding that requires a
bridge design to tolerate a given load and never mentions resonance has not
failed at verification when a resonant failure occurs; it failed earlier, at
the point where the obligation set was written. No amount of rigor at the
verification stage recovers a question the binding stage never asked.

\subsection{Compromised authorities}

A verification regime names a verifier, and naming a verifier is an act of
trust in whoever built, deployed, and maintains it. If that authority is
compromised, through insider action, supply-chain interference, or a captured
governance process, the verifier can return exactly the verdict its declared
procedure would produce on a correct implementation, while running a
different, compromised implementation the declared procedure never actually
checked. The check that a verdict was reached honestly is a separate concern
from the check that an obligation was satisfied, and this book returns to it
directly in the chapter on the threat model.

\subsection{Incomplete models}

Every obligation is stated in terms of some model of the domain it governs: a
model of financial risk, a model of structural load, a model of how a piece of
software will be used. A model can be internally coherent and still leave out
a variable that turns out to matter. Verification against an incomplete model
produces a verdict that is entirely correct with respect to the model and
entirely silent about the part of reality the model omitted.

\subsection{Unknown unknowns}

A closely related but distinct failure is the absence of any obligation at
all covering a class of risk nobody anticipated. An incomplete model at least
represents its domain with some variables missing. An unknown unknown is a
domain nobody thought to model in the first place. No obligation set can be
verified against a condition it was never written to check, and no
verification procedure can be faulted for failing to catch what it was never
told to look for.

\subsection{Valid proofs of the wrong proposition}

The most subtle failure mode is not an error in reasoning but a mismatch
between what was proven and what was needed. A verifier can construct a fully
valid derivation that obligation \(O_i\) is satisfied, where \(O_i\) itself,
though correctly proven, is not the proposition that mattered. This differs
from specification error in degree rather than kind: specification error is
usually a gap noticed only in hindsight, while a valid proof of the wrong
proposition can occur even when the specification looks adequate on its face,
because the gap lies in the translation from the informally intended
requirement to its formal encoding as \(O_i\). Interlude III returns to this
exact failure as a full worked case.

\section{Proof obligations and status}
\textbf{Limitation-setting rather than theorem-bearing.} This chapter's
purpose is to bound the scope of every sufficiency claim the rest of the book
makes, not to prove that absolute warrant is impossible in every formal system
whatsoever. That broader claim is not attempted here and would not be
established by the five cases above even if it were: each case is a
countermodel showing that a particular route to certainty fails in this
architecture's setting, not a general impossibility result about verification
as such. The chapters that follow are permitted to state exactly what a
discharged obligation guarantees; they are not permitted, and this chapter is
the place that forecloses it, to describe collapse, discharge, or a passing
verdict using the vocabulary of certainty. Chapter 14 revisits this same
territory once the formal verification apparatus exists and states the
narrower, precise non-claim illustrated by counterexamples there: a verifier's
own successful operation cannot certify the adequacy of the obligation set it
was checking against.

\chapter{The Witness Plane}
\section{Chapter thesis}

The natural first reading of a six-stage pipeline puts the ledger at the end:
an event is popped, refused or bound, transformed, verified, and finally, if it
survives, collapsed into history, with the ledger appearing afterward as a
record of what happened. That reading is wrong, and this chapter replaces it.
The ledger is not the last box. It is the substrate that witnesses every
boundary the pipeline crosses, including every boundary a candidate fails to
cross. A refused event, a failed transformation, an inconclusive verification,
and a stale collapse attempt all leave a mark on the ledger without ever
touching the system's operational history. The ledger's completeness does not
depend on anything reaching the end of the pipeline. It depends only on the
pipeline never taking a step it does not record.

This reframing matters for a concrete reason. If the ledger only appeared
after collapse, an observer would have no durable trace of anything the system
considered and declined. Every refusal, every failed proof, every stale
candidate would be as if it never occurred, and the distinctions built up
across Chapters 1 through 3---occurrence against endorsement, evidence against
warrant, sufficiency against certainty---would have nowhere to be recorded
once a decision went the other way. The witness plane is where those
distinctions become durable facts rather than momentary judgments.

\section{Two histories}

Two sequences are maintained side by side. The audit history, \(L_t\), is
extended by exactly one record at every step the system takes, regardless of
the step's outcome:
\[
  L_{t+1}=L_t\mathbin{\|}\ell_t.
\]
The operational history, \(H_t\), is extended only when a step succeeds at the
one transition authorized to extend it:
\[
  H_{t+1}\in\{H_t,\,H_t\mathbin{\|}\kappa_t\}.
\]

Which of the two alternatives obtains is not a matter of interpretation after
the fact. It is determined by the transition itself: every transition type
other than a successful collapse leaves \(H\) untouched, and a successful
collapse extends it by exactly one commit object \(\kappa_t\). Later chapters
define each transition precisely enough to state this as a proven property
rather than an assumption. The Early-Stage Non-Interference result of Chapters
6 through 10 covers the first five transition types; the collapse transition
itself, defined in Chapter 15, is the only one that can produce the second
alternative above.

Both sequences are append-only. Nothing already written to \(L_t\) or
\(H_t\) is later edited or removed. Redaction, treated fully in Chapter 19,
changes the availability of a payload a record refers to; it does not remove
the record itself or alter the sequence of hashes chaining prior records
together.

\section{The non-converse}

The relationship between the two histories runs in one direction only. Every
operational commitment must have a corresponding audit witness: each
\(\kappa\in H_t\) arose from some step, and that step wrote a record to
\(L_t\) documenting it. The converse does not hold. Most of what \(L_t\)
contains is not, and was never meant to become, an operational commitment. A
quarantined event, a repair request, a candidate that failed verification, a
stale collapse attempt rejected at the head check: each produces a durable
audit record, and none produces a \(\kappa\). An audit ledger that only ever
equaled the operational history in content would not be an audit ledger; it
would be a second copy of the same thing. The value of the witness plane is
precisely that it is larger than what it witnesses.

\section{Proof obligations and status}
\textbf{Load-bearing seed lemma.} This section proves, by induction on \(t\),
that every commit object in the operational history has exactly one
corresponding collapse-witness record in the audit history:
\[
 \forall\kappa\in H_t,\quad
 \exists!\,\ell^{\mathrm{collapse}}_\kappa\in L_t.
\]

The proof depends on five premises, each of which is either established
formally in a later chapter or is a property this book requires of any
implementation. They are stated here explicitly so the induction's scope is
visible before its steps are read.

\textbf{(P1) Append-only histories.} Neither \(L_t\) nor \(H_t\) ever loses or
overwrites a previously written entry.

\textbf{(P2) Sole writer.} Only the successful branch of the collapse
transition can extend \(H_t\), and it extends it by exactly one commit object
per successful collapse. This is proved for the five earlier transition types
in Chapters 6 through 10 and holds for collapse itself by the definition given
in Chapter 15.

\textbf{(P3) Atomic bidirectional publication.} When a collapse succeeds, the
append of \(\kappa_t\) to \(H_t\) and the append of its collapse-witness record
\(\ell^{\mathrm{collapse}}_{\kappa_t}\) to \(L_t\) occur as one indivisible
step. Neither occurs without the other, and no collapse-witness record is
created except through such a successful dual append. This final clause makes
explicit the no-orphan-witness fact needed by the uniqueness argument.

\textbf{(P4) Commit freshness.} Distinct successful collapses produce distinct
commit objects. This follows from the nonce-uniqueness clause of the collapse
guard, \(\neg\operatorname{Spent}(\eta)\), defined in Chapter 15, which
prevents any nonce, and therefore any commit object built from it, from being
produced twice.

\textbf{(P5) Typed records.} A ledger record of kind collapse-witness is keyed
to a specific commit object and cannot serve as a witness for any other commit
object. Records of other kinds---refusal records, failed transformation
records, failed verification records, and failed collapse records---are never
collapse-witness records for any \(\kappa\), regardless of their content.

\subsection{Base case}

At \(t=0\), \(H_0=\varnothing\). The claim holds vacuously: there is no
\(\kappa\in H_0\) for which a witness must be found.

\subsection{Inductive step}

Assume the claim holds at step \(t\): every \(\kappa\in H_t\) has exactly one
collapse-witness in \(L_t\). Consider step \(t+1\) and case split on the
transition performed.

\textit{Case 1: the transition is not a successful collapse.} By (P2),
\(H_{t+1}=H_t\). By append-only bookkeeping,
\(L_{t+1}=L_t\mathbin{\|}\ell_{t+1}\) for whatever record this step produces.
This record documents a pop, a refusal, a binding, a transformation, a
verification of any verdict, or a failed collapse attempt, and by (P5) none of
these record kinds is a collapse-witness for any commit object. So no existing
\(\kappa\in H_t\) gains a second witness, and since \(H_{t+1}=H_t\), no new
\(\kappa\) requiring a witness has appeared. Every
\(\kappa\in H_{t+1}\) still has exactly the one witness it had in
\(L_t\subseteq L_{t+1}\), and the claim holds at \(t+1\).

\textit{Case 2: the transition is a successful collapse producing
\(\kappa_{t+1}\).} By (P2) and (P3),
\(H_{t+1}=H_t\mathbin{\|}\kappa_{t+1}\) and, in the same step,
\(L_{t+1}=L_t\mathbin{\|}
\ell^{\mathrm{collapse}}_{\kappa_{t+1}}\). Two things must be checked: that
\(\kappa_{t+1}\) now has a witness and has only one, and that no previously
witnessed \(\kappa\in H_t\) loses its uniqueness in the process.

Existence for \(\kappa_{t+1}\) is immediate: its witness was just appended.
For uniqueness, suppose toward contradiction that \(\kappa_{t+1}\) already
had a witness in \(L_t\), before this step. By the no-orphan-witness clause of
(P3), the earlier witness could only have been created by an earlier successful
collapse that simultaneously appended the same commit object to \(H_t\). But
by (P4), \(\kappa_{t+1}\) is distinct from every commit object produced by an
earlier successful collapse. The supposed prior witness therefore cannot
exist, and exactly one is added now.

For every \(\kappa\in H_t\), the inductive hypothesis already supplies exactly
one witness in \(L_t\subseteq L_{t+1}\). The single record added at this step
is, by construction, keyed to \(\kappa_{t+1}\) specifically, and by (P5) it
cannot double as a witness for any distinct commit object. So no
\(\kappa\in H_t\) gains a second witness. Every element of
\(H_{t+1}=H_t\cup\{\kappa_{t+1}\}\) therefore has exactly one witness in
\(L_{t+1}\), and the claim holds at \(t+1\).

Both cases close the induction, establishing the result for all \(t\).

\subsection{Scope and forward reference}

This lemma is deliberately narrow. It covers only the successful-collapse path
under the five named premises, and it says nothing yet about crash recovery,
write-ahead durability, or implementations that persist \(H\) and \(L\)
through separate storage systems rather than one atomic operation. Chapter 18
strengthens exactly this result to cover recovery paths and a write-ahead
protocol in place of true atomicity, and should be read as extending this
proof's structure rather than replacing it. The seed lemma proved here is what
makes that extension possible: any later argument that ledger completeness
survives a crash must still reduce, in the crash-free case, to the induction
given above.

\interlude
\chapter*{Interlude I: A Request That Never Became an Act}

This interlude is not a chapter. It proves nothing and defines nothing
precisely. Its purpose is to walk one request through every stage this book
will eventually formalize, so that when the formal apparatus arrives beginning
in Chapter 5, it arrives as a description of something the reader has already
seen happen, rather than as machinery with nothing yet to grip.

\section*{The request}

A customer support agent notices that a customer was billed twice for the same
transaction. She submits a correction: reverse a duplicate charge of
\$412.17. Call this arriving event \(e_0\). It is observed at a definite time,
from a definite source, the agent's authenticated session, and it carries a
definite content: an account identifier, a signed amount, and a stated reason.

\section*{Refusal}

The system's intake policy, current version \(\sigma_t\), requires every
balance correction to reference the specific original transaction it corrects,
by its own identifier, not merely by a written description. Event \(e_0\)
does not carry that reference. The event is not silently dropped and it is not
silently accepted. It is refused, and the refusal is typed: this is a schema
deficiency, not a policy objection, a capability problem, or a question about
the agent's authority. The refusal decision records a disposition of repair, a
reason certificate naming the missing field, a stated requirement for what
resubmission needs to include, and a window in which resubmission remains
straightforward rather than requiring the request to start over. Event
\(e_0\) itself is preserved exactly as submitted. Nothing about the refusal
erases it; the refusal is itself a new fact added to the record, sitting
alongside \(e_0\) rather than replacing it.

\section*{Repair}

The agent looks up the original transaction, finds its identifier, and
resubmits. The new submission is not treated as though it had arrived out of
nowhere. It is a descendant: an event that names \(e_0\) as its parent and
supplies exactly the missing reference as its repair. Call this descendant
\(e_1\). The system checks \(e_1\) against the same requirement that
\(e_0\) failed, and this time the requirement is met. The original refusal of
\(e_0\) does not become false. It remains a true statement about \(e_0\),
evaluated under the policy in effect when \(e_0\) was submitted. What changes
is that a new, related event has now cleared the gate the first one did not.

\section*{Rebinding}

With \(e_1\) admitted, the system attaches the constraints that govern balance
corrections of this kind. Because the amount exceeds a threshold set by policy,
three obligations are generated before anything is computed or inspected: the
adjustment amount must match the discrepancy recorded against the referenced
transaction, a second approver distinct from the submitting agent must sign
off before the change takes effect, and the eventual audit trail must reference
the original transaction by its identifier. None of these obligations is
invented after the fact to justify a result. They are fixed now, as a contract
the eventual output will be judged against, whatever that output turns out to
be.

\section*{Transformation}

A kernel built for ledger adjustments computes the actual entry: the
account's new balance, given its old balance and the signed correction. The
computation is confined to exactly this task. It cannot reach beyond the
account in question, cannot consult unrelated records, and produces, alongside
its numeric result, a witness of how it arrived there. Nothing about this step
changes anything the system currently treats as authoritative. It produces a
candidate, not a commitment.

\section*{Failed verification}

The candidate is checked against the three obligations fixed at binding. The
amount matches the discrepancy: satisfied. The audit trail correctly
references the original transaction: satisfied. The second approver's
signature is present: it is not. The agent submitted alone. This third
obligation was never marked as one a policy officer could waive; it is required
outright for corrections of this size, precisely because requiring a second,
independent party to look at the change is the control this domain uses to
catch a class of error, or a class of abuse, that a single reviewer might miss
or might commit. The verification returns a verdict of failure, not because
anything was computed incorrectly, but because a specific, named, required
condition was not met. The system does not treat this as an error in its own
machinery, and it does not treat it as evidence that the correction is wrong.
It treats it as exactly what it is: one obligation, out of three, not yet
discharged.

\section*{Reopening}

The agent brings the correction to her manager, who reviews the transaction
reference and the computed amount and signs. A new event is submitted, again
not from nowhere: it names the failed attempt as its parent and supplies the
missing signature as its repair. The system processes this descendant through
binding and transformation again. The computation reproduces the same numeric
result it produced before, because nothing about the account or the correction
amount has changed; only the missing approval has been supplied. Verification
runs again, and this time all three obligations are satisfied. Nothing was
weakened to make this happen. The same requirement that blocked the first
attempt is checked again, in full, against the same standard, and this time it
passes because the actual condition it demands has actually been met.

\section*{Commitment}

A collapse proposal is formed: the verified candidate, its verification
result, a reference to the current state of the account ledger, the policy
under which this collapse is authorized, and a fresh, unused marker preventing
this exact proposal from ever being committed twice. The system checks that
nothing else has changed the account's ledger state since verification began,
checks that the proposal is authorized, checks that the candidate is indeed
collapsible, and checks that the posting mechanism is ready to accept the
entry. Every check passes. The correction is posted. The account balance
changes, and this change now belongs to what the system treats as settled and
current, not merely to what it has considered.

\section*{What remains afterward}

Two records exist afterward, and they are not the same record. The system's
operational history has grown by exactly one entry: the posted correction. The
system's audit history has grown by a great deal more: the original request,
its refusal, the repaired resubmission, the binding and its three obligations,
the first computed candidate, the failed verification naming exactly which
obligation was unmet, the second resubmission carrying the manager's approval,
the second computation, the successful verification, and finally the
commitment itself. An observer who only ever looks at the account's current
balance sees one clean change. An observer with access to the full record sees
that the change was neither immediate nor uncontested, that it was refused
once, repaired once, verified twice, and failed once in between, and that every
one of those steps is still there to be inspected, long after the balance
itself no longer shows any sign that the first attempt ever happened.

\section*{A preview of the graph}

Stated informally, ahead of the definitions to come, the identifiers this
story touched form a small lineage graph: \(e_0\), refused; \(e_1\), a repair
descendant of \(e_0\), admitted, bound, transformed, and found wanting at
verification; \(e_2\), a reopening descendant of \(e_1\), bound again,
transformed again to the same candidate, verified again and this time passed;
and finally a single commit object descending from \(e_2\)'s verified
candidate, the only object in this lineage that the operational history ever
came to contain. Chapter 5 begins supplying the formal definitions this graph
anticipates, starting with the event envelope itself.

\part{The Event Calculus}

\chapter{Events, Envelopes, and Identities}
\section{The event envelope}

Every distinction drawn informally in Part I needs a fixed object to attach to
once formal machinery begins. That object is the event envelope. Begin with an
identity-bearing event core
\[
  c(e)=\langle u,s,h,p\rangle,
\]
where \(u\) is a source-scoped occurrence token, \(s\) identifies the source,
\(h\) is a commitment to the content, and \(p\) is a stable reference to the
payload that content describes. The complete event envelope is
\[
  e=\langle id,u,t,s,h,p,m\rangle,
\]
where \(id\) is derived from the core, \(t\) is the observation time at which
the system first encountered the occurrence, and \(m\) is auxiliary metadata
that does not itself bear on identity or admissibility.

The occurrence token \(u\) distinguishes two intentionally distinct requests
that happen to have identical payloads, while remaining stable when a client
retries the same occurrence after losing an acknowledgement. It must be unique
within the source's identity domain and authenticated as part of the source's
submission. Observation time \(t\) cannot play this role because a retry is
observed at a different time from its first arrival.

Two further design choices are worth making explicit. First, the envelope
separates content commitment, \(h\), from payload reference, \(p\). The
commitment is what the rest of the architecture reasons about and incorporates
into later structures; the reference is how an authorized party locates the
actual bytes when it needs them. This separation is what later makes redaction
coherent in Chapter 19: a payload can become unavailable at its stable
reference without touching the commitment already incorporated into every
record built on top of it.

Second, \(t\) and \(m\) are deliberately excluded from the identity core.
Layout hints, display preferences, receipt timestamps, or free-text notes
belong in \(m\) precisely because they must never be capable of changing
whether two submissions denote the same occurrence. Placing something in
\(m\) is a declaration that it does not matter to the identity or admissibility
questions this machinery answers. If a field later proves relevant to either,
it was misclassified and must migrate into a versioned identity-bearing schema.

\section{Domain-separated identity}

An event's identifier is computed from its identity-bearing core, not assigned
and not computed from a structure already containing that identifier. Given a
canonical encoding function \(\operatorname{encode}\) and a fixed domain tag
\(\mathsf{domain}_e\) reserved for event cores specifically,
\[
 id(e)=\Hash\bigl(\mathsf{domain}_e\parallel
                   \operatorname{encode}(c(e))\bigr).
\]

The domain tag exists to prevent a coincidence that has nothing to do with
cryptographic weakness: two structurally different kinds of object, an event
core and a ledger record, for example, could encode to the same bytes if
nothing distinguished which kind of object was being hashed. Prefixing every
identifier construction in this book with a fixed, reserved tag specific to
the kind of object being identified prevents cross-domain equality by shared
encoding, independently of the collision resistance of \(\Hash\).

The encoding must be canonical, meaning that a given well-typed core has
exactly one encoding, not several equally valid ones. Raw concatenation without
canonical or length-prefixed encoding can make two different logical
structures produce the same byte string, the same failure mode as writing
\(ab\parallel c\) and \(a\parallel bc\) and expecting them to be
distinguishable without a delimiter. Section 5.4 makes precise what canonical
encoding is required to guarantee, and what it is not asked to guarantee.

\section{Equivalence and duplication}

Several notions of sameness apply to events, and failure to keep them apart
produces either false alarms or missed ones.

\textbf{Core equality} holds when two event cores encode to the identical byte
string. Given canonical encoding, this is the strongest notion tracked by
\(id(e)=id(e')\), subject to the collision-resistance assumption.

\textbf{Full-envelope equality} additionally requires equality of observation
time and auxiliary metadata. Two retry arrivals can therefore share an event
identifier without being the same observed envelope instance: the later
arrival has its own receipt time and may carry different non-identity metadata.

\textbf{Semantic equivalence} is broader and policy-relative. Two different
identity cores may express claims considered equivalent for a particular
downstream purpose. Such equivalence must not be silently built into the
identity hash. It belongs in an explicit normalization or equivalence relation,
because collapsing it into \(\operatorname{encode}\) would contradict the
injectivity required by the collision reduction below.

\textbf{Causal descent} is a relation between envelopes with different
identifiers, not a form of sameness. When an event is reopened, its descendant
carries new content, a fresh occurrence token, a new identifier, and an
explicit parent reference to the event it repairs. Chapter 8 treats this
relation formally. A descendant is related to its parent, but it is not the
same event.

\textbf{Retries} are arrivals that carry the same authenticated source token
\(u\), source \(s\), content commitment \(h\), and stable payload reference
\(p\). They therefore reproduce the same core and identifier even though their
receipt times differ. This is intended behavior: it allows intake to recognize
a resubmission as the same occurrence and deduplicate it instead of processing
the request twice. Reuse of \(u\) with a different \(h\) or \(p\) is not a
retry; it is a source-equivocation or malformed-submission event that must be
recorded and refused.

\textbf{Malicious replay} looks identical to a retry at the event-core layer:
an adversary resubmitting an authenticated core produces the same identifier a
legitimate retry would. The two cannot be distinguished by identity comparison
alone, and this chapter does not pretend otherwise. The danger concerns the
reuse of downstream authorization. Chapter 13 addresses it through freshness,
signature context, nonce uniqueness, and the current-head check at collapse.

\section{Proof obligations and status}

\textbf{Load-bearing reduction.} This section proves that a collision between
distinct identity-bearing event cores reduces to a collision in the underlying
hash function, under two named hypotheses, and states plainly what it has not
proven.

\textit{Hypothesis (Inj).} \(\operatorname{encode}\) is injective on the space
of well-typed identity cores:
\[
  c\neq c'\quad\Longrightarrow\quad
  \operatorname{encode}(c)\neq\operatorname{encode}(c').
\]

\textit{Hypothesis (Dom).} \(\mathsf{domain}_e\) is a fixed bit string of
known length, applied identically to every event-core encoding, and reserved
exclusively for this purpose elsewhere in the architecture.

\begin{proposition}[Event-identifier collision reduction]
Suppose \(c(e)\neq c(e')\) and \(id(e)=id(e')\). Then \(\Hash\) has a
collision.
\end{proposition}

\begin{proof}
By (Inj), \(c(e)\neq c(e')\) gives
\(\operatorname{encode}(c(e))\neq\operatorname{encode}(c(e'))\). Let
\[
 x=\mathsf{domain}_e\parallel\operatorname{encode}(c(e)),
 \qquad
 x'=\mathsf{domain}_e\parallel\operatorname{encode}(c(e')).
\]
Because \(\mathsf{domain}_e\) has fixed known length by (Dom), it can be
stripped from either input uniquely. Thus \(x=x'\) would force equality of the
two core encodings, a contradiction. Hence \(x\neq x'\). By definition,
\(id(e)=\Hash(x)\) and \(id(e')=\Hash(x')\). The assumed identifier equality
therefore gives \(\Hash(x)=\Hash(x')\) with \(x\neq x'\), which is a collision
for \(\Hash\).
\end{proof}

\subsection{What this does and does not establish}

This is a reduction, not a proof of collision resistance. It does not show
that \(\Hash\) has no collisions; it shows that an identifier collision
between distinct cores, given (Inj) and (Dom), cannot occur without one.
Collision resistance remains an imported cryptographic assumption, established
nowhere in this book and not attempted here \cite{katz-lindell}.

The two hypotheses are equally load-bearing and neither is a cryptographic
claim. Hypothesis (Inj) is a property of a serialization format, violated by
ambiguous field boundaries, non-canonical numeric formatting, or unordered-map
iteration. It applies to the identity core, not the full envelope: excluding
\(t\) and \(m\) from identity is an explicit projection, not a failure of
injectivity. Chapter 25's reference data model exists partly to specify a
format where (Inj) is checkable and enforced. Hypothesis (Dom) is an
engineering discipline maintained across the architecture: every identifier
construction introduced later, for ledger records, commit objects, and
collapse proposals alike, must use its own distinct, reserved tag. Appendix D
collects the complete registry of tags this requires.

A system satisfying (Inj) and (Dom) inherits exactly as much identifier
collision resistance as its hash function provides, no more and no less. A
system violating either hypothesis can produce identifier collisions with no
hash weakness at all, and no stronger hash function will repair a
non-canonical encoding or missing domain tag. Cryptographic strength and
specification discipline are independent sources of assurance, and a
reduction proof shows where their boundary falls.

\chapter{POP: Encounter Without Endorsement}
\section{Chapter thesis}

Before an event can be refused, bound, transformed, verified, or collapsed, it
must first be registered as something the system has encountered at all.
\(\Pop\) is the transition that does this registration, and nothing more. It
does not evaluate whether an encounter is welcome. It does not check the
encounter against policy. It converts raw, arriving material into the
well-typed envelope Chapter 5 defined, or determines that no such envelope can
be constructed, and in either case it writes a durable record that the
encounter happened. Everything this book calls admission---epistemic,
computational, operational, or irreversible---happens strictly after \(\Pop\),
never during it.

\section{Outputs and failure modes}

The transition \(\Pop\) takes raw arriving material, call it
\(e_{\mathrm{raw}}\), such as bytes on a wire, a message on a queue, or a
submitted form, and attempts to produce a well-typed envelope
\[
 e=\langle id,u,t,s,h,p,m\rangle
\]
as defined in Chapter 5. Five distinct tasks make up this attempt.

\subsection{Normalization}

Normalization converts \(e_{\mathrm{raw}}\) into the shape an envelope
requires: extracting a claimed source and occurrence token, separating payload
from descriptive metadata, and producing something canonically encodable.
Normalization can fail outright, before any of the remaining tasks begin, if
\(e_{\mathrm{raw}}\) cannot be parsed into that shape at all.

\subsection{Timestamping}

Timestamping assigns \(t\) as the time of observation by the system itself,
not a time claimed by whatever submitted \(e_{\mathrm{raw}}\). This detail is
not incidental. A client-supplied timestamp is a claim the client makes about
itself, and later chapters that reason about freshness and staleness need
\(t\) to be a fact the system asserts about its own observation, not a fact it
merely repeats on someone else's authority. For a recognized retry, the event
retains the first-observation time in its registered envelope while the new
arrival witness records the retry's own receipt time.

\subsection{Source authentication}

Source authentication verifies that \(s\) and its occurrence token \(u\)
correspond to a principal or session the system can authenticate, rather than
strings \(e_{\mathrm{raw}}\) merely claims. An unauthenticated or unverifiable
source produces intake failure prior to any question of whether that source is
authorized to perform a requested operation. Authorization is a separate
question addressed by later policy and binding stages.

\subsection{Payload commitment}

Payload commitment computes \(h\) from the payload's actual bytes and
establishes \(p\) as a stable reference through which an authorized party can
later retrieve them. This is where the separation argued for in Chapter 5,
between content commitment and payload reference, is first put into practice:
\(h\) is computed once, at intake, and everything built on top of the event
refers to the commitment rather than incorporating the payload bytes directly.

\subsection{Duplicate detection}

Duplicate detection checks whether the authenticated identity core
\(c(e)=\langle u,s,h,p\rangle\) has already been registered. A genuine retry
reproduces the same core and therefore the same \(id(e)\), allowing the event
registry to avoid creating a second independent event. The arrival itself is
not erased or ignored: each retry still produces a new POP witness in the
audit ledger, recording when it arrived and that it was deduplicated against
the previously registered event. Reuse of \(u\) with a different payload
commitment or reference is source equivocation rather than duplication and is
recorded as such.

\subsection{Intake failure}

Intake failure is the outcome when normalization, timestamping, source
authentication, or payload commitment cannot be completed at all, not because
the content is unwelcome, but because no well-typed envelope can be constructed
from what arrived. This differs in kind from the typed refusals Chapter 7
defines. Refusal presumes an envelope already exists and is being evaluated
against policy. Intake failure means the system never got that far. The
malformed network request of Chapter 1, missing a coherent content-length
field, fails here, at \(\Pop\), not at \(\Refuse\), precisely because there is
no well-formed \(e\) yet for a policy to evaluate.

\section{The first ledger witness}

Every one of these outcomes, successful registration or outright intake
failure, must leave a durable trace, and the two cases require different
minimum witnesses.

When normalization succeeds, the ledger records
\[
 \ell_t^{\mathrm{pop}}
 =\langle id(e),t,s,h,\mathit{dedup\_status}\rangle:
\]
enough to establish that this identified event was encountered at this time,
from this source, with this content commitment, and whether it was recognized
as a retry of something already seen. An implementation may distinguish the
event's first-observation time from the particular arrival time in an expanded
record.

When intake fails, there is no \(id(e)\) to record, because there is no
well-typed \(e\). The minimum sufficient witness commits to whatever raw bytes
can safely be represented, records the arrival time and whatever source
information could be determined even if it could not be authenticated, and
states which preceding task failed and why. Sensitive or adversarial raw input
need not be retained in plaintext; a commitment, bounded diagnostic sample,
and protected payload reference may be the appropriate witness.

This is the floor beneath the entire architecture: even when a system cannot
construct a single well-typed event from what it received, it must still be
able to say later that something arrived, roughly when, and roughly what
defeated the attempt to make sense of it. Without this floor, malformed input,
garbled transmission, and deliberately unparseable payloads would leave no
trace at all, precisely the kind of gap Chapter 22's threat model returns to
when considering adversaries who benefit from a system's inability to say what
it encountered.

\section{Proof obligations and status}

\textbf{Load-bearing non-interference lemma.} This section states and proves,
for the first time in this book, the Early-Stage Non-Interference Lemma that
Chapters 7, 9, and 10 instantiate for their transitions without repeating the
proof. The canonical statement is registered in Appendix C; it is proved here
because \(\Pop\) is the first transition defined formally, and the proof
technique is easiest to see in the simplest case.

\begin{lemma}[Early-Stage Non-Interference]
Let \(\tau\) be a transition whose complete defining semantics is given by a
finite set of equations specifying its outputs and effects on machine state.
Suppose \(H\) is either absent from those equations or appears only as a value
read by \(\tau\), and never as a state component assigned by \(\tau\). Then
for every \(t\), applying \(\tau\) at step \(t\) yields
\[
 H_{t+1}=H_t.
\]
\end{lemma}

\begin{proof}
The abstract machine's operational semantics, given fully in Appendix B,
adopts the standard frame convention that every state component not assigned
by a transition persists unchanged across that transition. By hypothesis,
\(H\) is never assigned by \(\tau\). Applying the frame convention to \(H\)
therefore gives \(H_{t+1}=H_t\) directly.
\end{proof}

The proof is intentionally almost immediate. Its content is not in the
derivation, which directly applies a convention fixed once in Appendix B, but
in the discipline the lemma imposes on every transition definition that
follows. Establishing the hypothesis requires an explicit check that \(H\)
never appears on the written side of the transition's equations. This check is
worth making precisely because it is easy to make silently false in an
implementation: a shortcut that caches a computed value into the same store
backing \(H\), for example, violates the hypothesis even if the local code
appears innocuous.

\subsection{Instantiation for \texorpdfstring{\(\Pop\)}{POP}}

The defining semantics of \(\Pop\) produce two kinds of result: a well-typed
envelope together with a POP witness, or an intake failure together with a
minimal salvage witness. In neither case does a transition equation assign a
value to \(H\). The lemma's hypothesis is satisfied by inspection, and it
follows that
\[
 H_{t+1}=H_t
\]
for every application of \(\Pop\), successful or failed. The transition
registers an encounter. It cannot, by the shape of its own definition, do
anything more than that.

\chapter{REFUSE: Typed Non-Admission}
\section{The refusal transition}

Once \(\Pop\) has produced a well-typed envelope, that envelope must still
clear a second gate before anything is bound, computed, or committed on its
behalf. The REFUSE gate evaluates the envelope against the policy version
current at decision time and returns one of two typed branches:
\[
 \Refuse_{\sigma_t}(e,H_t)\longrightarrow
 \begin{cases}
   \mathsf{ADMIT}(e,\alpha),\\
   \mathsf{REFUSAL}(d),
 \end{cases}
\]
where \(\alpha\) is a screening certificate authorizing the envelope to
proceed to BIND and
\[
 d=\langle r,\rho,\Delta,\tau,\pi\rangle
\]
is the typed refusal decision. Within \(d\), \(r\) is the disposition given to
the event, \(\rho\) is a typed reason certificate, \(\Delta\) states what a
resubmission would need to supply, \(\tau\) gives the conditions under which
the matter can be reconsidered, and \(\pi\) is the evidence object supporting
\(\rho\).

This is a decision about admissibility, not a judgment about the underlying
claim's truth or worth. A refusal states that this envelope, evaluated against
this policy and history at this decision time, does not clear the gate. It does
not state that the claim it carries is false, harmful, or without merit. The
admitted branch is equally bounded: it says that the event cleared screening,
not that its later transformation or commitment is warranted.

For a refusal, the corresponding ledger record is
\[
\begin{aligned}
\ell_t^{\mathrm{ref}}=\langle&id(e),h_{\mathrm{env}}(e),t_{\mathrm{obs}},
t_{\mathrm{dec}},\sigma_t,r,\rho,\Delta,\tau,\\
&h(\pi),\operatorname{parent}(e),\operatorname{sig}_G\rangle,
\end{aligned}
\]
where \(h_{\mathrm{env}}(e)\) commits to the complete registered envelope and
is distinct from its payload commitment \(h\). The record distinguishes
observation time from decision time, commits to the evidence object rather than
embedding it, tracks repair lineage through \(\operatorname{parent}(e)\), and
authenticates the deciding gate through \(\operatorname{sig}_G\). The admitted
branch analogously writes \(\ell_t^{\mathrm{admit}}\), committing to
\(id(e)\), \(h_{\mathrm{env}}(e)\), \(\sigma_t\), \(\alpha\), and the gate
signature. Thus clearing the gate is witnessed no less than failing it.

\section{Disposition algebra}

The refusal disposition \(r\) is drawn from a closed set, and its six values
are not interchangeable labels for the same act.

\textbf{Quarantine} holds an event pending resolution that is the system's
responsibility, not the submitter's, most commonly while an internal
investigation is underway. Nothing is asked of the submitter, and no repair
path is offered, because the obstacle is not a defect in what they submitted.

\textbf{Repair} applies when the obstacle is a specific, correctable defect in
the event itself, exactly the case Interlude I traced. The repair description
\(\Delta\) names what a descendant submission must supply.

\textbf{Evidence-only} retention preserves an event without pursuing a path
toward admission and without rendering a judgment on its eventual merit. This
is the disposition suited to the scientifically unsettled claim of Chapter 1:
retained because it may matter later, not admitted because nothing yet warrants
that.

\textbf{Deferral} differs from evidence-only retention by attaching an
explicit reconsideration condition in \(\tau\): a date, policy version, or
external event. Evidence-only retention makes no such temporal commitment.

\textbf{Redaction}, as a refusal disposition specifically, applies when an
event must be refused and its payload must additionally become unavailable,
usually for privacy or legal reasons attached to the content. This is narrower
than Chapter 19's general redaction mechanism, which may apply anywhere in a
record's lifecycle. Here, payload availability is removed at refusal, before
the event becomes a candidate for further processing.

\textbf{Denial} applies when evaluation is complete and final under current
policy and evidence: not merely unready, as in deferral, but considered and
rejected. A denied event remains fully present in the audit history. Denial
changes what happens next; it does not erase that the claim was made and
evaluated.

\subsection{Why purge is not an ordinary disposition}

No value of \(r\) authorizes deleting the fact that an event occurred. This
follows from the append-only architecture of Chapter 4. Removing an encounter
record would violate the completeness the witness plane exists to guarantee
and erase the trail that lets a refusal be checked later. Even redaction, the
closest disposition to what an ordinary system might call deletion, removes
only payload availability. The event identifier, content commitment, refusal
decision, and authorization for payload removal remain. What common usage
calls purging is represented here as authorized payload redaction under a
retained record, never removal of the authorization record itself.

\section{Reason certificates}

Six grounds populate \(\rho\), and distinguishing them matters because each
implies a different remedy, or the absence of one.

\textbf{Identity} grounds concern whether the authenticated principal is bound
to the identity or role claimed for this particular class of event, a finer
question than the bare session authentication \(\Pop\) performed.

\textbf{Schema} grounds concern whether content carried in a well-typed event
envelope satisfies the more specific structural requirements of current
policy, the missing transaction reference of Interlude I being the paradigm.

\textbf{Provenance} grounds concern whether the content's origin and chain of
custody meet the required standard, independently of whether the envelope is
well formed and its immediate source authenticated.

\textbf{Policy} grounds concern substantive prohibition: content that is well
formed, properly sourced, and adequately provenanced but falls outside what
current policy permits, such as an amount exceeding a hard cap.

\textbf{Capability} grounds concern possession of a particular delegated and
potentially revocable authority. A principal can be exactly who it claims to
be while lacking the capability token this event class requires.

\textbf{Temporal} grounds concern timing: arrival after a deadline, before a
waiting period has elapsed, or outside another declared validity window.

\subsection{Why \texorpdfstring{\(h(\pi)\)}{h(pi)} is a commitment, not a proof}

The evidence object \(\pi\) can be large, sensitive, or both: a fraud report,
excerpted logs, or an internal review record. The ledger commits only to
\(h(\pi)\) rather than embedding \(\pi\). Within an authenticated, timestamped
witness record, this secures a specific property: the evidence object used at
decision time is fixed against silent later substitution. It does not establish
that \(\pi\) actually supports \(\rho\), and hashing alone does not supply
authorship or non-repudiation. Those depend on the gate signature, provenance,
and declared evidence-evaluation procedure.

A committed evidence object and a sound reason certificate are separate
achievements. The first is a cryptographic fixation claim; the second requires
inspection of \(\pi\) against a declared standard for adequate support.
Treating \(h(\pi)\) as though it already certified \(\rho\) would repeat the
error Chapter 3 warned against: mistaking a mechanism securing one property
for a mechanism securing a stronger property it was never built to provide.

\section{Refusal as information}

A structured refusal certificate is more useful to a legitimate submitter than
a bare rejection. Naming the missing field, violated cap, or expired window is
what makes \(\Delta\) actionable rather than a guess. Yet the same specificity
can help an adversary map the system's defenses, learning which control caught
an attempt and adjusting the next one to route around it.

This tension has no uniform resolution across all six grounds. Schema and
temporal refusals are usually safe to disclose in detail. Identity, provenance,
and capability refusals sit closer to active defense, and repeated attempts
against these grounds are themselves evidence. Submitter-facing disclosure may
therefore be coarsened to report only a security-relevant refusal while the
full internal certificate remains available to authorized auditors.

What must not vary is the evidence commitment and full internal witness. The
submitter-facing message and internally retained record may differ in
granularity; the internal record must not be less complete than the actual
basis of decision. Chapter 18's ledger projections generalize this principle:
different audiences receive different views of one witness graph without
requiring the underlying graph to be impoverished.

\section{Proof obligations and status}

\textbf{Load-bearing non-interference lemma.} This section instantiates the
Early-Stage Non-Interference Lemma proved in Chapter 6. The generic lemma is not
reproved; only its hypothesis is checked.

The REFUSE gate produces either an admitted result and
\(\ell_t^{\mathrm{admit}}\), or a refusal decision \(d\) and
\(\ell_t^{\mathrm{ref}}\). No admitted certificate and no disposition in the
closed refusal set---quarantine, repair, evidence-only retention, deferral,
redaction, or denial---is defined by an equation assigning a value to \(H\).
Every branch writes to \(L\) alone. The hypothesis of Early-Stage
Non-Interference is therefore satisfied by inspection, and
\[
 H_{t+1}=H_t
\]
holds for every application of the gate. Whether the event clears screening or
receives a typed refusal, the decision extends audit history and leaves
operational history exactly as it found it.

\chapter{Reopening Without Revisionism}
\section{Descendants rather than mutation}

A refused event is never repaired by editing it. It is repaired by submitting
a new event that names the refused one as its parent and supplies what the
refusal's repair requirement requested:
\[
 e'=\operatorname{REOPEN}_{\nu_R}\bigl(id(e),\Delta',w,u'\bigr),
\]
where \(\Delta'\) is the claimed repair, \(w\) witnesses satisfaction of the
reconsideration condition \(\tau\), \(u'\) is a fresh source-scoped occurrence
token, and \(\nu_R\) versions the reopening specification. That specification
declares how \(\Delta'\) combines with the original payload to construct the
descendant's payload. The new event carries
\[
 \operatorname{parent}(e')=id(e).
\]
Once registered, \(e'\) is an ordinary event under Chapter 5, subject to
\(\Pop\), \(\Refuse\), and every downstream transition. Its only special
structural feature is its authenticated parent reference.

One degenerate case deserves separation. If no content has changed and the
submitter seeks only a new decision under a later policy or history, creating a
repair descendant would misdescribe what happened. The appropriate operation
is a reevaluation request over the same event identifier, producing a new
decision record under a new decision context rather than a new event. This
book does not develop \(\operatorname{REEVALUATE}\) in detail, but names it so
that REOPEN is not mistaken for the only way a policy change can affect a
prior refusal.

\section{Historical truth}

The refusal decision \(d\) made about \(e\) is a statement that the gate
produced a particular disposition when evaluating \(e\) under policy
\(\sigma_t\), history \(H_t\), and the evidence available at decision time.
Nothing that happens to a descendant changes that historical fact. If \(e'\)
is eventually admitted, verified, and collapsed, this does not by itself imply
that the original refusal was mistaken. It may instead show that \(e'\)
supplied something \(e\) lacked.

Append-only preservation must not be confused with correctness, however. A
later audit may show that the original gate misapplied \(\sigma_t\), relied on
bad evidence, or generated an erroneous reason certificate. Such a finding is
recorded as a later judgment about the earlier decision; it does not rewrite
the earlier record. The architecture preserves what was decided, by whom, on
what evidence, and under which policy, not a presumption that every preserved
decision was justified.

This distinction has two consequences. First, accountability evaluates the
decision against the policy, history, and evidence actually applicable at its
decision time, rather than silently substituting later conditions. Second, it
forecloses editing an old decision to make the record appear more consistent
in hindsight. Chapter 20 generalizes this as the invariant against in-place
revision of historical decisions.

\section{Repair graphs}

Treat every registered event identifier as a node and every authenticated
parent reference as a directed edge from parent to child. The resulting
structure, rooted at the original event, is a repair graph. It is often a tree
in practice but need not be a chain: a submitter may attempt several repairs of
the same parent, producing multiple children.

\subsection{Bounded retry}

Bounded retry is a policy layered over the graph rather than a structural
property. Policy may cap how many descendants of a root can be reopened before
mandatory escalation. This is the bounded-repair hypothesis used by Chapter
21's conditional liveness theorem; without it, the graph can grow indefinitely
and no termination argument follows.

\subsection{Acyclic repair}

Each successful REOPEN operation must name an already registered parent and
must introduce a fresh occurrence token \(u'\). Registration assigns every
new event a monotonically increasing creation rank
\(\operatorname{rank}(e)\). Therefore every parent edge satisfies
\[
 \operatorname{rank}(e)<\operatorname{rank}(e').
\]
This temporal-order invariant, rather than an appeal to payload equality,
forecloses cycles. Reuse of an existing occurrence token is handled as a retry
or equivocation by POP and cannot create a new lineage node or edge.

\subsection{Lineage queries}

Given an event identifier, its ancestor chain is recovered by following parent
references backward. Recovering all descendants requires an index over audit
records whose parent field names the identifier, because references point
backward only. Both directions matter: an auditor asks how many attempts
preceded success, while a reviewer asks what refusal a particular repair
answered.

\section{Proof obligations and status}

\textbf{Supporting lineage results.} These propositions are not central safety
theorems, but later audit and liveness arguments depend upon them.

\begin{proposition}[Fresh Descendant]
For
\(e'=\operatorname{REOPEN}_{\nu_R}(id(e),\Delta',w,u')\), where \(u'\) is
fresh in the authenticated source domain, \(id(e')\neq id(e)\), except with
the collision probability inherited from \(\Hash\).
\end{proposition}

\begin{proof}
The identity cores differ in at least the occurrence-token field, since \(u'\)
is fresh. Under (Inj) and (Dom) from Chapter 5, equality of the resulting event
identifiers would therefore induce a collision in \(\Hash\) by the
event-identifier collision reduction.
\end{proof}

\begin{proposition}[Non-Mutation]
No application of REOPEN alters or removes the refusal witness originally
written for its parent.
\end{proposition}

\begin{proof}
By premise (P1) of Chapter 4, existing ledger records are never overwritten.
REOPEN constructs a new event with a fresh identifier and writes only records
for that descendant. No defining equation reassigns a record keyed to the
parent. The original refusal witness therefore persists unchanged, although a
later audit may append a new record criticizing or superseding its judgment.
\end{proof}

\begin{proposition}[Acyclicity of Repair Graphs]
No repair graph whose parent references satisfy the registration-order
invariant contains a directed cycle.
\end{proposition}

\begin{proof}
Assume a directed cycle
\(e_0\to e_1\to\cdots\to e_k=e_0\). Every parent-to-child edge strictly
increases creation rank, so
\[
 \operatorname{rank}(e_0)<\operatorname{rank}(e_1)<\cdots
 <\operatorname{rank}(e_k)=\operatorname{rank}(e_0),
\]
which is impossible by irreflexivity of \(<\). Hence no directed cycle exists.
\end{proof}

\chapter{BIND: Constraint as Contract}
\section{The binding transition}

An admitted event does not proceed directly to computation. It first passes
through \(\Bind\), which consumes the admission certificate \(\alpha\) issued
by the REFUSE gate and attaches the requirements against which a candidate
output will eventually be judged. For fixed binding specification \(\nu_B\),
\[
 \Bind_{\nu_B}(e,\alpha,C,\Gamma_t)\longrightarrow
 \begin{cases}
  \mathsf{BOUND}(b),\\
  \mathsf{DEFERRED}(d_B),\\
  \mathsf{CONFLICT}(f_B),
 \end{cases}
\]
where \(C\) is the finite constraint set drawn from governing policies,
\(\Gamma_t\) is the binding environment describing which constraints apply and
how they relate, and
\[
 b=\langle e_C,\alpha,\beta,\mathcal O,\nu_B\rangle.
\]
Here \(e_C\) is the event paired with the constraints actually applied,
\(\beta\) certifies which \(C\), \(\Gamma_t\), precedence rules, and policy
versions produced the binding, and \(\mathcal O\) is the resulting obligation
set. The outcomes \(d_B\) and \(f_B\) witness incomplete binding context and a
detected constraint conflict respectively. They are outcomes of the total
binding function, not obligation sets disguised as successful bindings.

The architectural point is timing. The obligations are fixed before
\(\Transform\) produces anything and before \(\Verify\) inspects anything.
Interlude I's three obligations were not invented after a candidate correction
existed. They were fixed when the admitted request was bound, and the candidate
was judged against those terms. BIND commits the system, in writing, to what
will count as sufficient before it knows what computation will produce.

\section{Obligation generation}

Each member of \(\mathcal O\) is a four-part record:
\[
 O_i=\langle P_i,m_i,w_i,\delta_i\rangle,
\]
where \(P_i\) is the condition a candidate must satisfy, \(m_i\) is the
declared verification method, \(w_i\) is the class of witness required to
support a verdict, and \(\delta_i\) is the discharge policy stating whether and
under what authority the obligation may be waived rather than strictly
satisfied.

All four components are fixed from \((e,\alpha,C,\Gamma_t,\nu_B)\) alone.
None may be adjusted after a candidate exists. An obligation whose terms can
shift after output inspection is not a contract; it is post-hoc
rationalization written in contractual vocabulary.

\section{Constraint conflicts}

Four distinct problems can arise while turning constraints into obligations,
and each requires defined handling rather than ad hoc judgment at binding time.

\subsection{Unsatisfiable obligation sets}

Policies written in isolation can generate requirements no candidate could
jointly satisfy. BIND should catch syntactically obvious contradictions, such
as disjoint numeric intervals or a predicate paired with its direct negation,
and return \(\mathsf{CONFLICT}(f_B)\) rather than manufacture a contract known
to be hopeless. It cannot detect every unsatisfiable combination expressible
in a sufficiently rich constraint language; Section 9.4 records that limit.

\subsection{Policy precedence}

Where policies impose differing but compatible requirements on one topic,
\(\Gamma_t\) and \(\nu_B\) declare a fixed order of precedence. The rule may
select the more restrictive requirement or follow a jurisdictional hierarchy,
but \(\beta\) must record which rule was applied so an auditor need not infer an
unrecorded judgment.

\subsection{Local versus inherited constraints}

An event may carry locally declared constraints alongside constraints inherited
from policy. The default is asymmetric: a local constraint may tighten the
policy floor but may not loosen it. Otherwise a submitter could negotiate down
terms already fixed by governing policy. Domains permitting a different merge
must encode it explicitly in \(\nu_B\).

\subsection{Binding under incomplete information}

The environment \(\Gamma_t\) may lack context necessary for binding. BIND must
not guess and proceed as if it were present. It either applies a conservative
template that has been proven monotone under every permitted completion of the
missing information, or returns \(\mathsf{DEFERRED}(d_B)\). Merely calling a
template ``most restrictive'' is insufficient when candidate constraints are
incomparable; monotonicity must be a property of the binding specification.

\section{Proof obligations and status}

\textbf{Load-bearing well-definedness result.} The total result concerns the
binding outcome; obligation determinism concerns its successful branch.

\begin{proposition}[Binding-Outcome Well-Definedness]
Fix \(\nu_B\). For every finite, well-typed
\((e,\alpha,C,\Gamma_t)\), the mapping computed by \(\Bind_{\nu_B}\) returns
exactly one element of
\[
 \mathsf{BindingOutcome}
 =\mathsf{BOUND}(b)+\mathsf{DEFERRED}(d_B)+\mathsf{CONFLICT}(f_B),
\]
and is deterministic. Conditional on the \(\mathsf{BOUND}\) branch,
\(\mathcal O\) is uniquely determined.
\end{proposition}

\begin{proof}
Four properties of \(\nu_B\) establish the result.

\textit{Canonical ordering.} The specification fixes a total order on finite
constraint contributions, independent of collection or transmission order.

\textit{Fixed resolution.} It fixes a deterministic resolution function and
a canonical fold order for every pair of contribution forms admitted by its
domain. A detected syntactic contradiction maps to the conflict branch rather
than an undefined state.

\textit{Deterministic obligation identifiers.} Each resolved obligation is
assigned
\[
 id(O_i)=\Hash\bigl(\mathsf{domain}_O\parallel
                     \operatorname{encode}(O_i)\bigr),
\]
under the (Inj) and (Dom) hypotheses of Chapter 5. Identical resolved
obligations therefore receive identical identifiers.

\textit{Total handling of missing context.} When necessary information is
absent, the specification either selects a declared monotone conservative
template or returns the deferral branch. Thus no well-typed input leaves the
function stuck.

Canonical ordering and fixed resolution yield one resolved sequence. The
identifier rule yields one identifier per resolved obligation. Conflict and
missing-context cases map to distinct declared branches. Exactly one binding
outcome therefore results, and any successful obligation set is unique.
\end{proof}

\subsection{Explicit non-claim}

The proposition does not establish that a produced \(\mathcal O\) is jointly
satisfiable. General satisfiability for sufficiently expressive predicate
languages, including languages capable of encoding unbounded arithmetic or
arbitrary computation, is undecidable. BIND may and should decide declared
fragments and catch syntactically evident contradictions. Beyond them, it has
no general procedure for ruling out an impossible contract before TRANSFORM
and VERIFY attempt to discharge it. Deterministic generation is not adequacy,
correctness, or satisfiability.

\subsection{Non-interference}

The three BIND outcomes write binding, deferral, or conflict witnesses to
\(L\). None assigns \(H\). By the Early-Stage Non-Interference Lemma of Chapter
6,
\[
 H_{t+1}=H_t
\]
for every BIND application. Producing a contract, discovering a conflict, or
deferring for missing context does not commit an operational transition.

\interlude
\chapter*{Interlude II: Refusal Is Not Erasure}

Chapters 7 and 8 gave refusal a formal shape: a typed disposition, a reason
certificate, a repair path, and a lineage that survives whatever happens to a
repaired descendant. None of that machinery is original to computing. Every
institution that has had to decide what to accept and what to decline has
faced the same problem, and the oldest of them solved it in the same general
way this book does: by keeping a record of what was declined that can be every
bit as durable as the record of what was accepted.

\section*{What archives already knew}

An archive does not define its holdings by what it destroyed. A museum's list
of declined acquisitions, a library's record of what it chose not to shelve,
and a publisher's file of manuscripts it passed on are themselves historical
documents, sometimes more revealing of an era's priorities and blind spots
than the collection that resulted from them. Archival practice long ago
settled the distinction this book calls epistemic admission against
operational admission, without using those terms: something can be worth
remembering precisely because it was not selected, since the act of declining
it is itself a fact about the institution that declined it.

An archive that kept records only of what it accepted would not merely be a
smaller archive. It would be a different kind of thing, one unable to answer
what it once considered, which standards it applied, which alternatives it
excluded, or why it said no. The computational witness plane inherits this
institutional insight while making its mechanics explicit. Refusal changes an
event's admissibility; it does not retroactively prevent the encounter from
having occurred.

\part{Execution and Warrant}

\chapter{TRANSFORM: Computation Without Commitment}
\section{The transformation boundary}

Once an event has been successfully bound, a kernel is finally permitted to
compute something from it. The transition is total over two branches:
\[
 \Transform_{K,\theta}(b,H_t)\longrightarrow
 \begin{cases}
   \mathsf{CANDIDATE}(x),\\
   \mathsf{FAILURE}(f_K),
 \end{cases}
\]
where \(f_K\) is a typed timeout, resource, sandbox, dependency, or execution
failure. A successful candidate is
\[
 x=\langle y,\chi,T,\mathcal C,\nu_K,h(b),h(H_t)\rangle,
\]
where \(K\) is the kernel, \(\theta\) its configuration, \(y\) the candidate
result, \(\chi\) an execution witness, \(T\) the trace or a commitment to it,
\(\mathcal C\) the capabilities and resources consumed, and \(\nu_K\) the
kernel identity. The final two fields bind the candidate to the exact contract
and read-only operational context from which it was computed. Any reference to
\(H_t\) is read-only; neither branch may write operational history.

The whole of Chapters 1 and 3 converges on this transition's most important
property: it produces a candidate and nothing more. However correctly \(K\)
executes, \(x\) has no standing beyond being one computation's output until
VERIFY examines it against the obligations BIND fixed and COLLAPSE admits it.
Correct execution and authorization remain separate questions.

\section{Kernel identity}

The kernel identity is
\[
 \nu_K=\langle h(K),h(\theta),h(E),h(D)\rangle
\]
and commits separately to code, configuration, execution environment, and
dependencies. The separate fields prevent the phrase ``the same kernel'' from
hiding changes in parameters, runtime, libraries, hardware assumptions, data
files, or model weights. Each hash commits to a canonical manifest or artifact;
where a claim about the environment that actually executed is required, the
manifest must be joined to an attestation rather than treated as self-proving.

These commitments are also precisely what replay requires. A version label or
deployment tag is insufficient: replay must identify the exact code and the
declared context under which it ran.

\section{Capability confinement}

A kernel must be bounded along several dimensions before its output is worth
verifying.

\textbf{Authorization} determines which kernels and configurations may run
against which classes of bound event. Availability of a general-purpose kernel
does not license its use for every obligation class.

\textbf{Resource bounds} constrain execution time, memory, and other consumed
resources, producing a typed failure rather than indefinite execution or
unbounded consumption.

\textbf{Network access} is denied by default. When necessary, it must be
declared, and every response affecting computation must be captured or
committed as an input. Even nominally read-only network calls may produce
remote logging effects, so deployments requiring strict absence of external
effects must use recorded substitutes rather than live services.

\textbf{Clocks} are a common hidden source of nondeterminism. Every clock read
must be declared and its returned value captured. A client timestamp and a
system observation time remain distinct inputs.

\textbf{Randomness} must be deterministically seeded from a recorded seed or
treated as a declared nondeterministic channel whose actual outcome is
captured for replay.

\textbf{Side effects} that mutate shared or external state are prohibited
during TRANSFORM. Permitting them would let candidate production change the
world before verification or collapse.

\textbf{Other nondeterminism}, including scheduler interleavings, unordered
iteration, and floating-point reordering, must be eliminated or captured. An
undeclared channel removes the execution from the scope of replay equivalence.

\section{Execution witnesses}

Different obligation classes warrant different witness strengths. The chosen
style must be declared by \(\nu_K\) so VERIFY knows what evidence to expect.

A \textbf{full trace} records every relevant operation, offering high fidelity
at high storage and inspection cost. A \textbf{trace commitment} stores the
trace elsewhere and retains a content commitment and protected reference. A
\textbf{deterministic-replay witness} retains at least the complete input,
context, version, and nondeterminism manifest required to reproduce execution,
while omitting an operation-by-operation trace. \textbf{Proof-carrying
execution} supplies a machine-checkable proof about a declared computational
property, practical only for kernels and claims amenable to formalization.
An \textbf{attested environment} signs a statement that identified code ran
within an identified protected platform, shifting trust to the attestation
infrastructure and its trusted computing base.

No witness style is self-interpreting. Its verification relation, trust base,
and non-claims must be declared alongside it.

\section{Proof obligations and status}

\textbf{Load-bearing non-interference.} Both CANDIDATE and FAILURE append an
execution witness to \(L\); neither assigns \(H\). Early-Stage
Non-Interference therefore gives
\[
 H_{t+1}=H_t
\]
for every execution outcome.

\begin{theorem}[Conditional Transformation Replay Equivalence]
Assume: (a) \(K\) under \(\theta\) is deterministic in its bound input,
read-only history snapshot, initial state, and declared nondeterministic values;
(b) those inputs are canonically encoded; (c) \(\nu_K\) is identical across
original and replay executions; (d) \(h(b)\), \(h(H_t)\), and the readable
initial state are identical; (e) every clock value, random draw, network
response, scheduling choice, and other nondeterministic channel is eliminated
or replayed from its captured value; and (f) witness generation is itself
deterministic modulo a declared trace-equivalence relation \(\sim_T\). Then
replaying TRANSFORM produces the same \(y\) and an execution witness equivalent
under \(\sim_T\).
\end{theorem}

\begin{proof}
Under (a), the kernel is a mathematical function of the listed arguments. By
(b)--(e), each argument supplied to the replay is identical to its original
value. Function extensionality therefore yields the same \(y\). Under (f), the
witness mechanism receives the same semantic execution inputs and may differ
only in fields explicitly ignored by \(\sim_T\), so the replayed witness is
equivalent to the original under that relation.
\end{proof}

The hypotheses are substantive. If, for example, the kernel reads an
undeclared clock whose value influences rounding, two executions with the same
declared \(\nu_K\) may produce different \(y\). The theorem does not call such
an execution deterministic merely because it contains no explicit random-number
call. It excludes it because a semantic input was neither pinned nor replayed.

\chapter{VERIFY: The Discharge of Obligations}
\section{The verification transition}

A candidate produced by TRANSFORM is checked against the obligations BIND
fixed before it existed:
\[
 \Verify_{V,\sigma_v}(x,b,H_t)
 \longrightarrow
 v=\langle q,\mathbf d,\zeta,
          (\mathcal R_{\mathrm{block}},\mathcal R_{\mathrm{waived}}),
          \nu_V\rangle,
\]
where \(V\) is the verifier, \(\sigma_v\) its policy version, \(q\) the
overall verdict, \(\mathbf d\) the per-obligation discharge vector,
\(\zeta\) the verification witness, and \(\nu_V\) the verifier identity.
Like \(\nu_K\), \(\nu_V\) commits separately to code, configuration,
execution environment, and dependencies. A verifier is another program, not an
epistemically privileged observer, and Chapter 3's cautions apply to it in
full.

\section{Four verdicts}

A verification attempt terminates in exactly four reported ways.

\(\Pass\) holds when every obligation whose policy status is \(\Required\)
has been demonstrated \(\mathsf{SAT}\) under its declared method and witness
class.

\(\Fail\) holds when at least one required obligation is demonstrably
\(\mathsf{UNSAT}\), rather than merely unproven.

\(\Indeterminate\) holds when the verification procedure completed reliably
but the available evidence leaves at least one required obligation unresolved,
with no required obligation demonstrably violated.

\(\ErrorVerdict\) holds when the verifying procedure did not complete as a
valid run: a trusted wrapper detects a crash, timeout, resource exhaustion,
dependency failure, malformed result, or comparable execution fault. This is a
statement about the attempt, not about the true status of any obligation.

A binary pass/fail interface loses distinctions later stages need. Mapping
indeterminacy to failure conceals the difference between refutation and missing
evidence; mapping it to success admits the unproven. Mapping execution error to
either can make an infrastructure failure look like a substantive verdict.

\section{Discharge vectors}

Each obligation is recorded on two independent axes:
\[
\begin{aligned}
 d_i&=\langle s_i,a_i\rangle,\\
 s_i&\in\{\mathsf{SAT},\mathsf{UNSAT},\mathsf{UNKNOWN},
                 \mathsf{NOT\_CHECKED}\},\\
 a_i&\in\{\Required,\Waived\}.
\end{aligned}
\]

The status \(s_i\) is evidential: it records what the declared method shows
about \(P_i\). The status \(a_i\) is normative and procedural: it records
whether a shortfall is permitted, under the discharge policy \(\delta_i\)
fixed at binding, not to block commitment. A waiver changes operational
eligibility; it cannot change \(\mathsf{UNSAT}\), \(\mathsf{UNKNOWN}\), or
\(\mathsf{NOT\_CHECKED}\) into \(\mathsf{SAT}\).

Every waiver over a non-satisfied obligation cites the pre-encoded
\(\delta_i\), the current policy version granting the invoking authority, and
a witness that an authorized principal made this waiver for this obligation
instance. VERIFY may validate and report such an authorization; it may not
invent one after seeing the candidate.

\(\mathsf{NOT\_CHECKED}\) means the obligation was not established by a valid
completed verification run. In an otherwise reliable run, this may mean its
method was never executed. When the run ends in \(\ErrorVerdict\), normalization
sets every \(s_i\) to \(\mathsf{NOT\_CHECKED}\), even if partial internal
artifacts suggest that some checks completed before the fault. Those artifacts
may survive as quarantined diagnostics but are not reusable discharge results.
By contrast, \(\mathsf{UNKNOWN}\) means the declared method completed within a
valid run and produced genuinely inconclusive evidence.

The discharge vector is total over \(\mathcal O\): every obligation receives
one entry. The residual sets are
\[
 \mathcal R_{\mathrm{block}}
 =\{O_i\in\mathcal O\mid a_i=\Required
       \land s_i\neq\mathsf{SAT}\},
\]
\[
 \mathcal R_{\mathrm{waived}}
 =\{O_i\in\mathcal O\mid a_i=\Waived
       \land s_i\neq\mathsf{SAT}\}.
\]
The first blocks collapse. The second does not block it, but remains a required
disclosure containing the evidential deficit and its waiver authority.

\section{What a verifier verifies}

A verifier's report is recognized only if it satisfies a declared,
independently checkable relation:
\[
 \mathsf{Check}_{\nu_V}
 \bigl(h(b),h(x),\mathcal O,\mathbf d,\zeta\bigr)=1.
\]

Given commitments to the binding and candidate, the obligation set, discharge
vector, and witness, an authorized independent checker can determine whether
the report conforms to \(\nu_V\)'s declared relation. This does not prove that
\(\nu_V\) embodies an adequate standard, only that the supplied witness
supports the report under that standard. A vector lacking evidence sufficient
to evaluate this relation is an unsupported claim, not a verification result.

\section{Proof obligations and status}

\subsection{Verdict computation}

A trusted execution wrapper first determines whether a well-formed verifier
run completed within its declared resource and dependency conditions. If not,
\(q=\ErrorVerdict\). Otherwise, if any required obligation is
\(\mathsf{UNSAT}\), \(q=\Fail\). Otherwise, if every required obligation is
\(\mathsf{SAT}\), \(q=\Pass\). Otherwise, \(q=\Indeterminate\).

\begin{proposition}[Verdict Partition]
Conditional on the wrapper returning exactly one completion classification,
the four verdicts are jointly exhaustive and pairwise exclusive over every
verification attempt.
\end{proposition}

\begin{proof}
The decision procedure is a sequence of three ordered dichotomies: valid run or
execution error; given a valid run, some required demonstrated violation or
none; given none, complete required satisfaction or not. These cases exhaust
the possible classified outcomes. Because evaluation stops at the first
matching branch in a fixed order, no attempt reaches two branches.
\end{proof}

The proposition does not claim that every semantic fault inside \(V\) is
detectable. A compromised verifier may return a well-formed but wrong result,
and a compromised wrapper may misclassify the run. Those are trust and
correlated-failure questions for Chapters 12 and 22, not contradictions of the
partition over classified outcomes.

\begin{proposition}[Fault Containment]
\[
 q=\ErrorVerdict\quad\Longrightarrow\quad
 \forall i,\ s_i=\mathsf{NOT\_CHECKED}.
\]
\end{proposition}

\begin{proof}
This follows from the output-normalization rule enforced by the trusted
wrapper: on the error branch it constructs a fresh total vector indexed by
\(\mathcal O\) and assigns \(\mathsf{NOT\_CHECKED}\) to every evidential
component. It does not forward \(V\)'s partial vector. Thus no errored run can
export a reusable \(\mathsf{SAT}\) status.
\end{proof}

This is a reporting and containment property, not the claim that an earlier
partial check was logically false. It prevents later processes from salvaging
apparently good fragments of a run whose overall integrity failed. Full
re-verification from a clean attempt is required.

\subsection{Non-interference}

VERIFY writes its verdict, vectors, residual sets, and witness to \(L\). It
does not assign \(H\). By Early-Stage Non-Interference,
\[
 H_{t+1}=H_t
\]
for every verdict, including \(\Pass\). Passing verification makes a candidate
eligible for the collapse guard; it does not itself commit the candidate.

\chapter{Verifier Independence and Correlated Failure}
\section{Chapter thesis}

Placing verification in a separate stage from transformation is not, by
itself, enough to make it independent. Independence is a property of what the
two stages share, not of how many boxes separate them. If VERIFY parses output
with the same faulty library TRANSFORM used to produce it, both may agree
because they share the blind spot that made the result look correct. If the
obligation set encodes a common misreading of an ambiguous requirement,
separate implementations do not catch the gap. If one authority approves,
deploys, or maintains both kernel and verifier, one compromise may place a
matching flaw in both. This chapter names these shared failure channels because
a stage that shares them is not doing the independent work its position in the
pipeline appears to promise.

\section{Diversity and quorum}

Several practices strengthen verification only against particular failures.

\textbf{Heterogeneous verifiers}, built by separate teams with different
languages and libraries, reduce common implementation faults. They do nothing
against a specification error every team encoded identically.

\textbf{Threshold acceptance} is only as strong as the independence of the
verifiers counted toward its quorum. A quorum sharing one critical dependency
creates apparent redundancy without protection against that dependency.

\textbf{Disagreement records} matter independently of resolution. Divergent
verdicts may reveal ambiguous obligations, model differences, or unequal
evidence. Majority resolution should not erase that diagnostic fact.

\textbf{Majority voting over correlated implementations} can create positive
miscalibration. Three matching reports are stronger than one only to the
extent that their relevant failure modes are genuinely separate. Counting one
shared interpretation or dependency three times does not make it independent.

\section{Recursive verification}

Requiring every verifier to be verified invites a regress: its verifier would
itself require checking. This architecture truncates the regress explicitly at
a named bounded trust root, such as attested hardware, a minimal reference
implementation, a formally verified kernel, or a designated institution.

The truncation must be narrow and accompanied by an assurance statement. A
hardware root may be trusted for a vendor-attested execution identity while
providing no independent assurance of internal logic. A reference verifier may
have a machine-checked soundness proof whose force still depends on the proof
assistant and the specification formalized. The goal is not to eliminate the
regress but to expose where internal verification ends and imported trust
begins.

\section{Proof obligations and status}

\textbf{Load-bearing negative result.} Let dependency \(D\), whether a
library, specification, model, or authority, be shared by verifiers
\(V_1,\ldots,V_k\). Suppose \(D\) contains a fault for which there is a class
of candidates \(X\) such that every verifier routing obligation \(O_i\)
through \(D\) misreports \(s_i\) on every candidate in \(X\). Let
\(p_X=\Pr[x\in X\mid D\text{ has the stated fault}]\).

\begin{proposition}[Correlated Verifier Failure Lower Bound]
Conditional on \(D\) having the stated fault,
\[
 \Pr[V_1,\ldots,V_k\text{ all misreport }s_i\mid D\text{ faulty}]
 \geq p_X,
\]
independently of \(k\). If the fault is present with probability \(p_F\) and
the candidate distribution conditional on its presence has mass \(p_X\), then
the unconditional joint-failure probability is at least \(p_Fp_X\), absent
additional dependence information.
\end{proposition}

\begin{proof}
Conditional on the stated fault, the event \(x\in X\) implies that all
\(k\) verifiers misreport. It is therefore a subset of the joint-failure event,
so the latter has probability at least \(p_X\). The expression contains no
factor depending on \(k\). For the unconditional form, the intersection of
fault presence and \(x\in X\) implies joint failure; multiplying its conditional
probability by \(p_F\) yields \(p_Fp_X\) under the stated conditional model.
\end{proof}

\subsection{Why naive multiplication fails}

If each verifier has observed marginal error rate \(\varepsilon\), an
independence model estimates unanimous error as \(\varepsilon^k\). That
calculation is valid only when the relevant error events are independent. If a
substantial component of \(\varepsilon\) arises from \(D\), the joint error is
bounded below by the shared-fault term and need not shrink as verifiers are
added. The arithmetic is not wrong; the independence model is.

\subsection{Diversity that addresses the dependency}

Independent reimplementation with genuinely separate libraries weakens
library-level correlation. Independently derived formalizations of the
informal requirement address specification-level correlation. Separate
organizational authority and supply chains weaken common compromise channels.
Diverse hardware and runtimes expose environment-specific faults. None of
these substitutes automatically for another: hardware diversity does not
repair identical source logic, and separate codebases do not repair a shared
misreading of one specification.

Surface diversity does not suffice. Teams using different languages may still
bind to the same load-bearing dependency, consult the same ambiguous source,
or answer to the same authority. Assurance claims must therefore inventory
shared dependencies and justify independence at the level of the failure being
mitigated rather than infer it from organizational or syntactic appearances.

\chapter{Freshness, Substitution, and Replay}
\section{Candidate identity}

A verification result is meaningful only if tied unambiguously to the exact
candidate and binding it verified. Chapter 10 placed \(h(b)\) inside the
candidate as \(x.h_b\), and the signed verification result commits to both
\(h(x)\) and that binding commitment. Define
\[
 \operatorname{SameCandidate}(x,v)
 \;:\!\iff\;
 h(x)=v.h_x\ \land\ x.h_b=v.h_b.
\]
An adversary cannot pair a legitimate passing verification for \(x\) with a
different \(x'\) unless it finds a distinct object with the committed hash or
alters the authenticated verification result.

\section{Freshness}

Five freshness predicates apply at collapse, and each protects a different
time-of-check to time-of-use boundary.

\textbf{Elapsed-time freshness} bounds the interval between verification and
collapse. A verified market quote may become useless without any change to the
candidate's internal content.

\textbf{State-head freshness} requires the history head embedded in \(x\) and
the collapse proposal to equal the actual current head at the atomic collapse
operation.

\textbf{Policy freshness} requires the policy version used for binding,
verification, and waiver authority to remain current or explicitly
grandfathered.

\textbf{Capability freshness} requires every authority used by verification
or collapse to remain unrevoked at collapse time.

\textbf{Evidence freshness} checks the separate validity windows of external
evidence, which may expire before the verification result's general time limit.

The predicate \(\operatorname{Fresh}(v)\) is the conjunction of the freshness
requirements declared by the binding and collapse policy. Their reference
versions, validity intervals, and registry identifiers must be committed into
the authenticated verification and collapse records so the predicate can be
re-evaluated rather than trusted as an opaque Boolean.

\section{Attack catalogue}

\textbf{Certificate replay} resubmits an already successful proposal.
\textbf{Candidate substitution} pairs a passing result with an unverified
candidate. \textbf{Stale authorization} relies on a revoked capability or
superseded policy. \textbf{Rollback} presents an old history head as current.
\textbf{Equivocation} presents different current heads to different parties,
defeating the premise of freshness rather than its local equality check.
\textbf{Time-of-check to time-of-use gaps} name the general interval in which
state, policy, evidence, or authority can change after verification and before
collapse.

\section{Proof obligations and status}

\textbf{Load-bearing conditional security theorem.} Let \(\mathcal A\) be a
probabilistic polynomial-time adversary able to observe \(L\), delay, reorder,
and replay messages, substitute candidates, and corrupt every principal except
the trusted collapse arbiter, verifier-signing authority, freshness registries,
and clock named below.

Assume: (i) \(\Hash\) is collision-resistant; (ii) signatures authenticating
\(v\) and \(c\) are existentially unforgeable under chosen-message attack;
(iii) spent-nonce tracking is linearizable and rejects every reuse of \(\eta\);
(iv) the current-head comparison is an atomic compare-and-append performed by
an uncorrupted arbiter against the true head; and (v) elapsed-time, policy,
capability, and evidence freshness are evaluated fail-closed against an
authenticated consistency snapshot of the required registries and a trusted
clock, with all checked versions and validity windows committed into the signed
proposal.

\begin{theorem}[Freshness and Substitution Resistance]
Under (i)--(v), the probability that \(\mathcal A\) causes the collapse guard
to accept a replayed, substituted, or stale candidate is negligible in the
security parameter, except for the explicitly assumed failure probabilities of
the trusted arbiter, registries, clock, and nonce store.
\end{theorem}

\begin{proof}
Suppose \(\mathcal A\) succeeds with non-negligible probability. Since the
guard is a conjunction, success must evade at least one relevant check.

\textit{Candidate substitution.} Presenting \(x'\neq x\) while satisfying
\(h(x')=v.h_x=h(x)\) yields a hash collision, unless the adversary changes the
committed value in \(v\), which requires a signature forgery. The same argument
applies to substituting a different binding commitment.

\textit{Certificate replay.} Reusing a successful proposal reuses its nonce.
By (iii), the atomic \(\neg\operatorname{Spent}(\eta)\) check rejects it. A
random nonce collision is also rejected; nonce-space size concerns availability
and accidental rejection, not acceptance of a spent proposal. Altering \(\eta\)
inside the authenticated proposal requires a signature forgery.

\textit{State-head staleness or rollback.} By (iv), a proposal bound to an old
head fails the atomic current-head comparison. Presenting two distinct
histories with one head commitment requires a collision in \(\Hash\).

\textit{Other freshness failures.} By (v), expired evidence, elapsed time,
revoked capability, or obsolete ungrandfathered policy causes
\(\operatorname{Fresh}(v)\) to fail. Changing a committed version or validity
window requires a signature forgery or hash collision. Acceptance despite an
authentic stale registry value requires failure of the trusted freshness
infrastructure, which the theorem records as an external assumption rather
than bounding cryptographically.

A union bound over collision, forgery, and the explicitly imported trusted
component failures yields negligible adversarial success under the stated
hypotheses.
\end{proof}

\subsection{What the theorem does not cover}

The atomic, non-equivocating arbiter and authenticated freshness infrastructure
are indispensable. A dishonest arbiter may present different heads to
different parties without finding a collision, forging a signature, or
reusing a nonce. Likewise, a compromised capability registry can report a
revoked authority as current while every local cryptographic check succeeds.
Chapter 16 names the linearizability assumption required to prevent forks;
Chapter 22 places these trusted components inside the threat model. Stronger
hashing cannot repair a false source of current state.

\chapter{Verification Is Not Truth}
\section{Chapter thesis}

A verifier can execute correctly, suffer no dependency failure, avoid every
correlated-implementation problem Chapter 12 names, satisfy the hypotheses of
Chapter 13's security theorem, and still certify an obligation set that fails
to capture what mattered. This is not the claim that verification is
unreliable. Even a maximally reliable verifier cannot exceed what it was told
to check. Chapter 3 promised to return to this point once the formal apparatus
could state it narrowly. With that apparatus now available, the limitation is
almost definitional.

\section{Specification horizons}

Three questions sit at increasing distance from what VERIFY answers.

\textbf{Internal validity} asks whether the procedure correctly assessed the
declared predicates using their methods and witness classes. Chapters 11--13
specify the conditional assurance available here, including verdict semantics,
correlated-failure limits, candidate identity, and freshness.

\textbf{Operational adequacy} asks whether \(\mathcal O\), taken as a whole,
captures what the domain requires for commitment. This concerns the content of
the obligation set supplied by BIND. VERIFY takes that set as an input and has
no semantic rule allowing it to infer completeness relative to an unstated
domain requirement.

\textbf{World-relative success} asks whether acting on a candidate satisfying
an adequate obligation set produces the desired real-world outcome. Even a
well-specified and correctly verified contract can fail here because the model
omitted a relevant variable or the world contains an unknown condition.

Consider an autonomous braking system whose obligations require stopping
distance to remain within a bound for a declared friction coefficient and
speed. The verifier executes flawlessly and the trial evidence supports
\(q=\Pass\). Internal validity holds. Suppose, however, that the binding has no
predicate for a sharp localized change in friction during the stop. VERIFY did
not fail; it answered the question it received. The obligation set omitted the
question that later mattered, and execution reliability cannot manufacture a
predicate that BIND never supplied.

\section{Revisable assurance}

An incident, near miss, or new research can revise confidence in the adequacy
of the earlier obligation set. The domain community may add the missing
predicate so later bindings do not repeat the gap. This does not rewrite the
historical record that the earlier verifier produced \(\Pass\) for the narrower
set under the policy then in effect. Under this chapter's premise of a correct
verification run, the report was accurate about what it checked. What hindsight
changes is the judgment of whether that was the right thing to check.

This is the historical-preservation principle of Chapter 8, not a rule that
old judgments must forever be considered correct. If later evidence shows the
earlier run itself was faulty, that finding is appended as a new judgment about
the preserved result rather than used to edit the original record.

\section{Proof obligations and status}

\textbf{Documented non-claim with a definitional proposition.}

\begin{proposition}[What \texorpdfstring{\(\Pass\)}{PASS} Certifies]
By the semantics of \(\Pass\) alone, \(q=\Pass\) for \((b,x,v)\) certifies
that every obligation whose policy status is \(\Required\) received
\(s_i=\mathsf{SAT}\) under its declared method \(m_i\), supported by evidence
of the declared witness class \(w_i\), as reported in \(\mathbf d\) and
checkable through \(\mathsf{Check}_{\nu_V}\). It assigns no additional
architectural meaning concerning the completeness or adequacy of
\(\mathcal O\), or the real-world consequences of collapse.
\end{proposition}

\begin{proof}
Chapter 11 defines \(q=\Pass\) as the reliable-run branch in which every
required obligation is \(\mathsf{SAT}\). The definition of
\(\mathsf{Check}_{\nu_V}\) ranges over the committed binding, candidate,
obligations, discharge vector, and witness. Neither definition quantifies over
unstated domain requirements, the fidelity of an informal-to-formal
translation, or downstream consequences. Therefore no such claim belongs to
the semantics of \(\Pass\) itself. Stronger conclusions may follow when
separate domain premises are supplied, but their additional force comes from
those premises rather than from VERIFY alone.
\end{proof}

\subsection{Why this is not a G\"odelian result}

G\"odelian incompleteness concerns what sufficiently expressive formal systems
can prove in settings involving internalized syntax and self-reference. The
present limitation requires none of that machinery. It is a scope fact about
an interface: VERIFY evaluates the obligations supplied to it. The braking
counterexample shows what lies outside that declared input. Invoking
incompleteness would dress an ordinary semantic boundary in borrowed authority
while obscuring the more useful question of who constructed \(\mathcal O\),
from which model, and under what process of revision.

\interlude
\chapter*{Interlude III: The Proof That Proved Too Little}

Chapter 3 mentioned a binding that requires a bridge to tolerate a given load
and never mentions resonance. Chapter 14 gave that sentence its precise
content: a verifier can execute flawlessly and still certify an obligation set
that never asked the right question. This interlude gives the sentence a full
life, tracing one certification through the pipeline without a refusal or
failed verification, and showing what that cleanliness did and did not
guarantee.

\section*{The submission}

A structural engineering firm submits a certification request for a new
pedestrian bridge: lightweight, materially efficient, and intended for a river
crossing that has needed one for years. The request is well formed, arrives
from an authenticated and appropriately credentialed source, and POP registers
it without incident.

\section*{Admission}

REFUSE finds nothing to object to. Schema and provenance are complete, the firm
holds the required capability, and timing raises no temporal concern. The event
is admitted on its first attempt, with no repair descendant.

\section*{The obligations}

BIND draws from the certification code then in force. Static load capacity must
meet the design specification for foot traffic, material fatigue ratings must
satisfy the intended service life, and the overall safety factor must clear the
regulatory minimum. These are serious obligations. Their ordering is canonical,
precedence is unambiguous, and a second application of the same binding
specification produces the same set.

\section*{Computation and verification}

The kernel performs load-distribution calculations, material-stress modelling,
and fatigue projection. The candidate exceeds every declared threshold.
VERIFY checks each obligation. Static load, fatigue rating, and safety factor
are all \(\mathsf{SAT}\); \(q=\Pass\), with neither blocking nor waived
residuals. A reviewer inspecting exactly these obligations has nothing in the
discharge vector to flag.

\section*{Commitment}

The collapse guard finds the actor authorized, the candidate collapsible, the
head current, the nonce unspent, and the certification mechanism ready. The
certification becomes operationally authoritative. Construction proceeds and
the bridge opens.

\section*{What no obligation asked}

Months later, during a crowded crossing, the bridge begins to sway. The cause
is not static overload but pedestrian footsteps coupling with a low-damping
lateral mode. Pedestrians adjust their gait for balance, strengthening the
synchronization. The bridge closes for inspection; no one is seriously hurt,
but retrofit is lengthy and costly.

\section*{The audit that found no pipeline violation}

Review recovers the intake witness, admission certificate, binding and its
three obligations, execution witness, clean discharge vector, and collapse
record. Obligation generation was deterministic. Computation and verification
were correct under their declared models. Replay reproduces the same result.
Nothing was skipped or falsified.

The obligation set did not contain a predicate for pedestrian-synchronized
lateral response in a low-damping design. The phenomenon was known in the
broader literature but absent from the jurisdiction's adopted code. Whether
that absence should be attributed to negligence, institutional lag, or a
reasonable limitation of knowledge is a separate governance inquiry. At the
pipeline level, the defect was a specification gap rather than a failure to
discharge the specification supplied.

\section*{What changes and what does not}

The code is revised to require explicit pedestrian-synchronized lateral
analysis below a damping threshold. The revision does not edit the earlier
record. That record continues to show a design satisfying every obligation its
binding generated, verified without procedural error and committed under a
guard that found no declared defect. The bridge still swayed. Both facts remain
visible, allowing later reviewers to distinguish a failure of execution from a
failure in what the institution had learned to ask.

\part{Commitment and Effect}

\chapter{COLLAPSE: The Commitment Boundary}
\section{The collapse proposal}

Everything preceding this chapter exists to be checked at one final boundary.
A collapse proposal gathers what must be true for a verified candidate to enter
operational history:
\[
 c=\langle x,v,h(H_t),\sigma_c,\eta\rangle.
\]
Its identifier is derived without self-reference:
\[
 id_c=\Hash\bigl(\mathsf{domain}_c\parallel\operatorname{encode}(c)\bigr).
\]
The commit object is then
\[
 \kappa=\langle id_c,h(H_t),h(x),h(v),\sigma_c,
                   \eta,t_c,a_c\rangle.
\]
It records the proposal, prior head, candidate and verification commitments,
policy, nonce, commit time, and committing authority. The new head is derived
only afterward:
\[
 h(H_{t+1})=
 \Hash\!\left(\mathsf{domain}_H\parallel h(H_t)
 \parallel\operatorname{encode}(\kappa)\right),
\]
and belongs to the history header or append result, not to \(\kappa\) itself.
This avoids a circular pre-image in which the commit object contains the hash
whose computation already depends upon that object.

The compact operational object \(\kappa\) commits to \(v\) through \(h(v)\).
Its corresponding audit witness carries the fuller disclosure record required
for inspection, including any waived deficits. Operational history need not
duplicate every audit field in order to commit cryptographically to it.

\section{The collapse guard}

A proposal succeeds only if five independent checks hold:
\[
\begin{aligned}
G(c,H_t)={}&\operatorname{Authorized}(c)
 \land\operatorname{Collapsible}(x,v)\\
&\land(h(H_t)=c.h_H)
 \land\neg\operatorname{Spent}(\eta)\\
&\land\operatorname{EffectsReady}(x).
\end{aligned}
\]

Authorization concerns the committing actor, not merely the original
submitter. Head agreement enforces state freshness atomically with append.
Nonce uniqueness rejects replay. Effect readiness determines whether the
internal posting mechanism or Chapter 17's prepared external-effect mechanism
can proceed. Collapsibility is defined against the blocking residual set:
\[
\begin{aligned}
\operatorname{Collapsible}(x,v)\iff{}&q=\Pass
 \land\mathcal R_{\mathrm{block}}=\varnothing\\
&\land\operatorname{Fresh}(v)
 \land\operatorname{SameCandidate}(x,v).
\end{aligned}
\]

\section{Principal proposition}

\begin{proposition}[Collapse Guard Respects Waivers]
Two properties hold:

\textnormal{(a)} collapse eligibility depends on
\(\mathcal R_{\mathrm{block}}\), not on whether
\(\mathcal R_{\mathrm{waived}}\) is empty; and

\textnormal{(b)} every waived deficit is explicitly disclosed in the collapse
audit witness and cryptographically committed by \(\kappa\) through \(h(v)\).
\end{proposition}

\begin{proof}
For (a), the definition
\[
 \operatorname{Collapsible}(x,v)
 \Longleftrightarrow
 q=\Pass\land\mathcal R_{\mathrm{block}}=\varnothing
 \land\operatorname{Fresh}(v)
 \land\operatorname{SameCandidate}(x,v),
\]
contains no emptiness test for \(\mathcal R_{\mathrm{waived}}\). Holding the
displayed conjuncts fixed therefore fixes collapsibility regardless of the
waived set.

For (b), the successful-collapse rule constructs
\(\ell^{\mathrm{collapse}}_\kappa\) by copying, for every
\(O_i\in\mathcal R_{\mathrm{waived}}\), the obligation identifier, unchanged
\(s_i\), authorizing \(\delta_i\), waiver-policy version, authority, and witness
from \(v\). No transition rule rewrites these fields. The operational object
stores \(h(v)\), so alteration of the verification result or disclosed waiver
data breaks that commitment, subject to hash collision resistance. Thus waiver
changes eligibility without laundering a deficit into \(\mathsf{SAT}\), and
the deficit remains inspectable in the witness plane.
\end{proof}

\section{Typed collapse failure}

A false guard produces a typed failure rather than a generic rejection.

\textbf{Stale head} is detected during preflight when the proposal's base head
already differs from the current head. Remedy may require reverification or
rebinding, depending on what the intervening change affected.

\textbf{Replay} means \(\eta\) is already spent. Both malicious replay and an
innocent duplicate proposal are rejected. A new attempt requires a newly
authorized proposal with a fresh nonce.

\textbf{Unauthorized} means the committing actor lacks current standing under
\(\sigma_c\).

\textbf{Verification invalid} covers failure of a collapsibility conjunct:
non-PASS verdict, blocking residual, stale verification, or candidate mismatch.
The audit record retains the failed sub-conjunct rather than flattening these
reasons into one opaque code.

\textbf{Effect unavailable} means the declared realization mechanism is not
ready. This failure may be transient, but retry remains subject to freshness
and nonce rules.

\textbf{Commit conflict} means preflight observed a matching head, but another
proposal won the subsequent atomic compare-and-append. This is a genuine race
lost at the serialization point, distinct diagnostically from a proposal known
to be stale before attempting the atomic operation.

Every failure appends \(\ell^{\mathrm{collapse\_failure}}\) to \(L\) and
leaves \(H\) unchanged:
\[
 H_{t+1}=H_t,
 \qquad
 L_{t+1}=L_t\mathbin{\|}\ell^{\mathrm{collapse\_failure}}.
\]
Only the successful atomic branch appends \(\kappa\) to \(H\) and its matching
collapse witness to \(L\).

\chapter{Concurrency and the Moving Head}
\section{Compare and append}
The collapse guard's head-agreement check, \(h(H_t)=c.h_H\), is not a
convenience check performed before the real work of appending happens. It is
the same primitive that underlies every lock-free concurrent data structure:
compare a shared value against an expected value, and perform the update only
if they match, as a single, indivisible operation. If the comparison and the
append were instead two separate steps---check the head, then append---a window
would open between them during which another collapse could change the head.
The original attempt could then append on top of a head that was no longer
current by the time it acted. This is the time-of-check to time-of-use gap from
Chapter 13 applied to the moment operational history advances. Atomicity is
not a performance optimization. It is what makes the head-agreement check mean
what it claims to mean.

\section{Competing valid candidates}
Two candidates can each be fully and correctly verified while being jointly
inconsistent if both are committed. Consider two withdrawal requests against
the same account, submitted close together. Each is individually verified
against the balance as it stood when its verification ran, and each satisfies
every obligation in its own obligation set. If both could be committed, their
combined value might exceed the balance even though neither verification,
taken alone, could see the other candidate racing toward commitment.

The head-based atomic compare-and-append solves this problem not by making
verification smarter, but by serializing commitment. Whichever withdrawal
reaches the atomic append first changes the head. The second attempt, built
against the old head, then fails head agreement and must be reconsidered
against the new state. Under that state, its balance obligation may now fail,
correctly and for exactly the reason joint consistency requires: the first
commitment spent value on which the second candidate depended.

\section{Rebinding after conflict}
A head conflict establishes that some operational fact changed; it does not by
itself establish how far backward through the pipeline a candidate must go.
The remedy depends on which earlier artifact the intervening commit may have
invalidated.

If the new head is proven irrelevant to both the derivation of \(y\) and the
evaluation of every obligation, the existing candidate may be reverified
against the new head. This is the narrowest case: the history moved, but no
state read by TRANSFORM and no fact used by VERIFY changed.

If the intervening commit changed state read by TRANSFORM, while the binding
and obligation set remain current, reverification of the old candidate is not
enough. The candidate itself may be stale. The system must rerun TRANSFORM
against the new state and then VERIFY the newly produced candidate.

If the intervening change altered the governing policy, the binding
specification \(\nu_B\), or any constraint context from which \(\mathcal O\) was
derived, full rebinding is required. The system returns to BIND, generates a
new obligation set under the current specification, and carries the resulting
binding through TRANSFORM and VERIFY again. Reverifying against a stale
contract would silently hold a candidate to requirements the system no longer
considers authoritative.

\section{Distributed ordering}
A single sequencer, one designated arbiter processing collapse attempts in a
strict order, provides linearizability directly. It is the simplest design to
reason about and the clearest single point of failure.

Consensus protocols let multiple replicas agree on one commit order while
tolerating some replica failures. They can provide the same abstract
linearizability guarantee as a sequencer, at the cost of latency and protocol
complexity. Paxos and Raft are representative constructions
\cite{lamport-paxos,ongaro-raft}.

Causal partial orders relax the demand for one global total order. Commits
without causal dependence may proceed without being ordered against one
another. This changes what ``history never forks'' means: history becomes a
directed graph rather than a single chain. The theorem below assumes a total
order and does not transfer automatically to this setting.

Sharded histories give separate state partitions independent local heads and
arbiters. Commits confined to different shards need not coordinate. A collapse
that must update several shards atomically, however, requires machinery beyond
the theorem in this chapter.

Domain-local collapse occupies a middle position. Each domain maintains its
own history and arbiter, while explicitly declared cross-domain obligations
trigger synchronization. The common case remains local; cross-domain
consistency becomes a visible and deliberately more expensive exception.

\section{Principal theorem}
\textbf{Theorem (History Never Forks).} Assume:

\emph{(H1)} \(H_0\) is a unique, well-defined chain, whether the empty history
or a single agreed genesis state.

\emph{(H2)} Every application of COLLAPSE performs its head comparison and
append as one linearizable atomic compare-and-append operation. This operation
is enforced either by a single arbiter processing attempts serially or by a
consensus mechanism providing the equivalent guarantee that every observer
sees operations as if they occurred in one agreed total sequence.

Then, for every \(t\), \(H_t\) remains a unique chain, and at most one candidate
referencing a given \(h(H_t)\) as its expected head succeeds in extending it.

\emph{Proof.} Proceed by induction on \(t\). The base case holds by (H1). For
the inductive step, assume \(H_t\) is a unique chain, with one agreed value of
\(h(H_t)\). Consider all collapse attempts submitted with that expected head.
By (H2), these attempts take effect as if processed in one agreed serial order.
Failed attempts leave the head unchanged. The first attempt in that order that
satisfies every conjunct of the guard succeeds and atomically replaces the
shared head with \(h(H_t\mathbin{\|}\kappa)\). No later attempt expecting the
old head can then succeed, because its head-agreement conjunct is false at its
own linearization point. Thus at most one candidate referencing \(h(H_t)\)
extends the history.

If no attempt succeeds, \(H_{t+1}=H_t\). If one succeeds,
\(H_{t+1}=H_t\mathbin{\|}\kappa\) for the unique successful commit object.
Either result is one well-defined chain agreed by every observer under (H2),
so the induction closes. \(\blacksquare\)

\textbf{What the theorem does not get for free.} Neither Chapter 5's
identity-core construction nor Chapter 13's freshness and substitution guards
supplies (H2). Identity protects object distinction; freshness protects a
correctly functioning head comparison. Neither prevents two non-linearizable
replicas from independently accepting different commits against what each
believes is the current head. In that case, two internally valid chains can
diverge without any hash collision, signature forgery, or nonce reuse. This is
the equivocation attack Chapter 13 deferred. Preventing it requires an actual
single sequencer or a correctly implemented consensus or compare-and-swap
service. It is an engineered infrastructure premise, not a consequence of
cryptographic hygiene elsewhere in the architecture.

\chapter{Irreversible Effects}
\section{The external-world problem}
A ledger append is entirely within a system's own control. It touches storage
the system manages, and Chapter 4 established the conditions under which the
append and its audit witness can be published atomically. A physical
actuation, a wire transfer through an external network, or a shipment dispatched
through a third party is a different kind of operation. It reaches
infrastructure this architecture does not control, through a network or
physical link that can fail, delay, or duplicate messages. There is generally
no way to make “record this locally” and “cause this to happen out there”
one indivisible operation because they do not occur within one transactional
boundary.

Suppose a system calls an external wire-transfer API and the connection times
out before a response arrives. The request may never have reached the bank, or
the bank may have completed the transfer while only its response was lost.
Those possibilities are indistinguishable to the caller from the timeout
alone. This ambiguity is not cured by better exception handling. The protocol
below is designed to survive it without pretending it away.

\section{Prepare, actuate, finalize}
Three steps replace the impossible atomic operation:
\[
 \operatorname{PREPARE}(c)\to\iota,
 \qquad
 \operatorname{ACTUATE}(\iota)\to\epsilon,
 \qquad
 \operatorname{FINALIZE}(\iota,\epsilon)\to\kappa.
\]
PREPARE records a specific intent \(\iota\) durably before any request is
sent outside the system. The intent includes the collapse proposal, the
external operation to be requested, and an idempotency key derived from its
committed fields. A process resuming after failure can therefore discover that
an actuation may be pending without reconstructing intent from volatile state.

ACTUATE first writes a durable attempt witness keyed to \(id(\iota)\), then
uses \(\iota\)'s stored idempotency key in the external request. If a response
arrives, it returns a receipt \(\epsilon\), which records either the external
system's attestation of success or a typed failure. The attempt witness is
written before transmission because a crash after transmission but before
local bookkeeping would otherwise make a genuine retry look like a first
attempt. The receipt may still be lost even when the external effect succeeds.

FINALIZE combines \(\iota\) and a conclusive \(\epsilon\), obtained either
from the original response or from later reconciliation, into the commit object
\(\kappa\) appended to operational history. An inconclusive transport failure
is not a receipt of non-occurrence and cannot authorize finalization.

Between successful PREPARE and FINALIZE lies an inherent interval in which the
effect may already have occurred while \(H\) reflects neither success nor
failure. The prepared intent belongs to the machine's durable pending-intent
state \(P_t\), introduced in Appendix B, rather than to \(H_t\). If a deployment
treats \(P_t\) as authoritative operational state, it instantiates Chapter 2's
\(\mathrm{Op}\land\lnot\mathrm{Ir}\) case. If it reserves
\(\mathrm{Op}\) strictly for \(H\), PREPARE is instead a durable procedural
commitment below operational admission. Conversely, successful ACTUATE followed
by a crash before FINALIZE produces the orphaned-effect case
\(\mathrm{Ir}\land\lnot\mathrm{Op}\).

\section{Recovery semantics}
Exactly-once external execution cannot be guaranteed from this side of the
boundary. The defensible target is durable intent, at-least-once delivery where
appropriate, idempotent processing where the external system supports it, and
explicit reconciliation wherever ambiguity remains.

Durable intent is \(\iota\), retained before actuation. Every actuation attempt
uses the same stored idempotency key and produces a local attempt witness.
Receipts are retained and referenced by FINALIZE whether they arrive promptly
or through a later status query. Reconciliation examines a stalled intent,
queries the external system using that key, and either finalizes an effect
shown to have occurred, records a conclusive failure, or retries under the
external system's declared duplicate-handling contract.

When an effect occurred but later must be opposed, the remedy is a compensating
action: a refund, reversal, corrective shipment, or other new event. The
original consequence cannot be deleted from the world. Compensation therefore
resembles Chapter 8's descendant repair: it adds a causally related act rather
than rewriting the act that preceded it.

\section{Proof obligations and status}
\textbf{External assumption, not an internal theorem.}

\textbf{Proposition (Intent Determinism and Local Duplicate
Detectability).} Let \(k(\iota)\) be the idempotency key stored by PREPARE, and
let \(A(\iota)\) be the durable sequence of actuation-attempt witnesses for the
intent. Then: (a) every ACTUATE invocation for the same \(\iota\) presents the
same key; and (b), assuming the pending-intent and attempt stores are durable,
an invocation is locally identifiable as a retry exactly when
\(A(\iota)\neq\varnothing\) before its new attempt witness is appended.

\emph{Proof.} PREPARE computes \(k(\iota)\) once as a pure function of the
intent's committed fields and stores it inside the immutable intent. ACTUATE
reads this value rather than generating a new one, proving (a). By the
transition's required order, the first invocation finds no prior attempt
witness and durably appends one before transmission. Every later invocation
finds at least that witness, even if the first transmission or its response was
interrupted by a crash. Under the durability premise, the emptiness test
therefore distinguishes a first local attempt from a retry, proving (b).
\(\blacksquare\)

The attempt journal is necessary. The mere presence of \(\iota\) proves that
PREPARE occurred, not that ACTUATE was previously attempted; an intent can
remain prepared for an arbitrary period before its first transmission.

Neither part proves what an external system does when it receives the same key
twice. ACTUATE-idempotency is a capability contract of that integration, not a
theorem of Spherepop. Where the contract holds, the external system absorbs
duplicate requests. Where it does not, Spherepop can prove only that its own
duplicate attempts are visible and available to reconciliation. Even this
weaker guarantee matters: an unpreventable duplicate need not also be a silent
one, and a detected duplicate can trigger a separately recorded compensating
action.

\chapter{The Ledger as a Witness Structure}
\section{Ledger records}
Every audit record accumulated across the preceding chapters shares one common
shape:
\[
\begin{aligned}
 \ell_i=\langle&seq_i,kind_i,subject_i,h(payload_i),policy_i,\\
                &parents_i,t_i,actor_i,sig_i,h(\ell_{i-1})\rangle.
\end{aligned}
\]
The sequence number \(seq_i\) is assigned by the ledger and fixes append order
independently of any claimed wall-clock time. The typed field \(kind_i\)
distinguishes pop, refusal, reopening, binding, transformation, verification,
collapse, collapse failure, preparation, actuation attempt, reconciliation,
and related witnesses. It makes Chapter 4's typed-record premise mechanically
checkable. The subject identifies the event, binding, candidate, intent, or
commit object concerned. The payload hash commits to detailed content while
permitting that content to be stored under a separate access regime.

The remaining fields name the governing policy, causal parents, timestamp,
actor, and authenticating signature. Finally, \(h(\ell_{i-1})\) chains the
record to its predecessor. Altering a committed field changes the record hash,
which invalidates the successor's predecessor commitment and propagates toward
the anchored head under the assumptions stated below.

\section{Tamper evidence and its limits}
Tamper evidence does not establish correctness. A dishonest actor can write a
false record initially and leave it untouched thereafter; the chain preserves
that falsehood perfectly. It protects integrity after recording, not truth at
recording time.

Nor does it establish availability. A well-chained record may be lost,
destroyed, inaccessible, or deliberately withheld. Replication, backup, and
access governance remain separate requirements.

Hash chaining also does not establish non-equivocation. A ledger operator can
present different observers with different chains, each internally consistent.
Chapter 16's linearizability premise governs \(H\); it does not automatically
govern \(L\). If the two histories use different publication mechanisms, the
audit ledger can equivocate even while operational history remains unique.
Closing that gap requires an authoritative, consistently published ledger head
or an equivalent consensus or transparency mechanism.

Finally, the chain does not make \(t_i\) honest. A corrupted or skewed clock
can supply a false timestamp that the chain then preserves faithfully.
Independent timestamping, external anchoring, or attested clocks are needed
when temporal truth is itself load-bearing.

\section{Ledger projections}
Different audiences require different views of the same witness graph. An
audit view exposes structural fields and the signature chain while respecting
payload access restrictions. An operational view selects collapse witnesses.
A provenance view follows parent relations for one subject. A policy view
selects records governed by a particular policy version. A security view
selects relevant refusal, authentication, capability, and failure records. A
privacy view withholds unauthorized payload references while retaining the
structural fact that a record exists.

These are derived projections, not independently maintained histories. A
faithful projection must satisfy two conditions relative to its declared
filter: it fabricates no source record, and it omits no source record that the
filter requires. Selective disclosure may hide fields or payloads, but any
claim of completeness for a projection requires either access to the source
graph or a cryptographic inclusion-and-completeness proof supplied by the
implementation. The hash chain alone supplies neither.

\section{Principal theorems}
\textbf{Load-bearing ledger completeness.} Chapter 4 proved that every
\(\kappa\in H_t\) has exactly one collapse witness in \(L_t\) under atomic dual
publication. Real deployments may store \(H\) and \(L\) separately, so this
chapter admits a weaker implementation premise.

\emph{(P3') Recoverable write-ahead publication.} Before either visible history
is extended, a single durable transaction record \(w_\kappa\) commits to both
\(\kappa\) and \(\ell_\kappa^{\mathrm{collapse}}\). Each application to \(H\)
or \(L\) is keyed by \(id(w_\kappa)\) and is idempotent. Recovery replays every
transaction not marked complete on both sides, applying only the missing
append, before ordinary processing resumes.

\textbf{Theorem (Strengthened Ledger Completeness).} Under (P1), (P2), (P4),
and (P5) from Chapter 4, together with either atomic dual publication or (P3'),
\[
 \forall \kappa\in H_t,\qquad
 \exists!\,\ell_\kappa^{\mathrm{collapse}}\in L_t
\]
holds at every quiescent machine state, including every state exposed after
crash recovery has completed.

\emph{Proof.} The crash-free cases are Chapter 4's induction. Under (P3'), a
crash can occur before either append, after the \(H\) append only, after the
\(L\) append only, or after both. Recovery consults the durable transaction
record and the idempotent application markers. In the first case it applies
both sides; in either middle case it applies only the missing side; in the last
case it applies neither. Thus every transaction reaches one append on each side
and never two. Premises (P4) and (P5) preserve uniqueness across distinct
commits. Recovery completes before the resulting state is exposed, so every
visible \(\kappa\in H_t\) has exactly one matching witness in \(L_t\).
\(\blacksquare\)

The qualification “quiescent” is essential. During recovery, or while a
write-ahead transaction is durably prepared but only one side is visible, the
pointwise invariant can temporarily be false. The protocol guarantees
restoration before normal readers or writers are admitted, not metaphysical
simultaneity across separate stores.

\textbf{Load-bearing tamper evidence.}

\textbf{Theorem (Anchored Hash-Chain Tamper Evidence).} Assume canonical,
injective encodings for ledger records, domain separation, collision resistance
of \(\Hash\), and an authentic reference to a later head \(r_n=h(\ell_n)\).
If a committed field of \(\ell_i\), \(i\leq n\), is altered while the anchored
head remains fixed, validation of the chain segment from \(i\) through \(n\)
rejects except with negligible collision probability.

\emph{Proof.} Altering a committed field changes the canonical encoding of
\(\ell_i\). Equality of its old and new hashes would therefore yield a
collision by Chapter 5's reduction. If the successor remains unchanged, its
stored predecessor commitment no longer equals the recomputed hash. If an
adversary rewrites the successor to repair that link, the successor's own
encoding and hash change; the same dilemma recurs at every later record. To
arrive at the unchanged authentic head \(r_n\), the rewritten suffix must
eventually produce the old hash from a distinct input, again yielding a
collision. Hence validation against the authentic anchor rejects except with
negligible probability. \(\blacksquare\)

For the special counterfactual in which the suffix is recomputed honestly from
the altered record without adversarially compensating fields, every descendant
hash differs except with negligible probability. The operational detection
claim is narrower: a verifier must obtain and validate the intervening chain
segment, or receive an authenticated proof structure designed for abbreviated
verification. Possession of a head hash alone does not locate an edit and does
not eliminate the need to authenticate a path to it.

The theorem establishes none of correctness at initial entry, availability,
global non-equivocation, or honest timestamping. It establishes only that an
alteration cannot be reconciled with an already authentic later anchor without
breaking the hash assumption.

\chapter{Redaction, Secrecy, and Cryptographic Erasure}
\section{Persistent occurrence, removable content}
Chapter 5 separated an event's content commitment from the retrievability of
its payload. The event identifier, commitment, and ledger position persist;
what redaction can remove is the ability to resolve the immutable payload
handle \(p\) to plaintext. Redaction is therefore structurally closer to
sealing than expungement. It cannot make the encounter disappear without
abandoning append-only history.

The distinction is implemented through an external resolver
\(\mathsf{Resolve}_t(p)\), not by editing the event envelope. Before redaction,
the resolver may return ciphertext, a storage location, or authorized
plaintext. After redaction it returns a typed tombstone or access denial. The
handle \(p\) itself remains fixed inside the event's identity core.

\section{Mechanisms}
A bare hash of low-entropy content permits dictionary testing. A randomized
commitment
\[
 h=\Hash(\mathsf{domain}_{payload}\parallel
         \operatorname{encode}(payload,salt))
\]
resists this only while a high-entropy salt remains secret; if the salt was
published, deleting it later adds no hiding. Encryption with key destruction
leaves ciphertext in place while rendering it unreadable, conditional on
strong encryption and the absence of surviving key copies. A resolver
tombstone distinguishes deliberate redaction from a payload that never existed.
Access revocation is weaker but reversible because privileged recovery remains
possible.

Selective disclosure, commonly through committed substructures such as Merkle
trees, can reveal authorized fields with inclusion proofs. Zero-knowledge
proofs can establish a declared property of hidden content without revealing
the content itself. Neither mechanism proves that the hidden payload was true;
each proves only consistency with a commitment and satisfaction of the
specified relation.

\section{Governance}
Every redaction produces a typed authorization record naming its subject,
authorizer, policy provision, mechanism, time, and any review or reversal
conditions. The redaction record is itself append-only. A resolver may cease
returning the payload, but the fact that access was removed and the authority
under which it was removed cannot vanish with it.

\section{Proof obligations and status}
\textbf{Proposition (Redaction Preservation).} Let Chapter 5's immutable
identity core be
\[
 c(e)=\langle u,s,h,p\rangle,
 \qquad
 id(e)=\Hash(\mathsf{domain}_e\parallel\operatorname{encode}(c(e))).
\]
Suppose redaction changes only the resolver state associated with \(p\), or
destroys a key needed to open the referenced ciphertext, while leaving the
canonical envelope and \(c(e)\) unchanged. Then \(id(e)\), \(h\), and the
canonical envelope hash \(h(e)\) remain unchanged.

\emph{Proof.} None of the arguments to \(id(e)\) or \(h(e)\) changes.
Redaction changes the time-indexed resolution relation, not \(p\) itself.
Both functions therefore receive byte-identical canonical encodings before
and after redaction, while \(h\) is retained directly. \(\blacksquare\)

This formulation preserves Chapter 5 rather than revising it retroactively.
Replacing \(p\) inside the envelope with a tombstone would change an
identity-bearing field and therefore change both the event encoding and,
absent a collision, its identifier. The tombstone belongs in the resolver and
the redaction ledger record.

The result establishes commitment continuity, not later inspectability or
non-repudiation of erased plaintext. Once every opening, key, salt, and payload
copy is destroyed, the system may retain a commitment it can no longer open.
It can still truthfully say that a specific value was committed at a specific
time; it may no longer be able to demonstrate what that value was.

\interlude
\chapter*{Interlude IV: The World Does Not Roll Back}
Chapter 17 built a protocol around a fact that deserves plain expression:
some effects cannot be undone, only answered. A database rollback operates on
the same internal state as the original transaction. Once an effect enters the
world, restoration of a later state is not erasure of the path that produced
it.

\section*{Physical actuation}
A valve that fired has fired; material that was cut remains cut; a vehicle
that moved occupied the intervening trajectory. Closing the valve, replacing
the material, or returning the vehicle is a new act with new costs and risks.
The later action may restore a selected variable while leaving every collateral
consequence of the interval intact.

\section*{Public communication}
A published statement may be corrected, retracted, or removed from its
original host, but it cannot be unread, unremembered, or withdrawn from every
copy made before revocation. The correction often reaches fewer people than
the claim. Chapter 19's redaction controls future retrieval within governed
stores; it cannot recall information already transported beyond them.

\section*{Human decision-making}
Reliance makes institutional decisions path-dependent. An applicant resigns
from another position, a patient undergoes a procedure, or a student allows a
deadline to pass because an assurance was given. Rescission does not restore
the prior life trajectory. It adds a new decision atop consequences already
incurred.

This is why compensation must never be described as undo. Durable intent,
receipts, reconciliation, and corrective descendants are honest precisely
because they preserve the difference between opposing an effect and making it
never have happened. Part V now asks what safety and liveness can mean once
that difference is admitted.

\part{System Properties and Failure}

\chapter{Safety Invariants}
\section{Chapter thesis}
This chapter proves nothing new. It collects the preceding results, identifies
their dependencies, and distinguishes abstract-machine invariants from
conditional infrastructure claims.

\section{Six invariants}
No unbound transformation and no verification without a candidate hold by the
typed transition signatures of Chapters 10 and 11. These are unconditional
inside the abstract machine, though an implementation still owes a refinement
argument showing that its interfaces enforce those types.

No collapse without a passing certificate follows from Chapter 15's guard:
\(q=\Pass\), candidate identity and freshness, and an empty blocking residual
set are necessary conjuncts. No operational commitment without exactly one
audit witness is Chapter 4's seed result strengthened by Chapter 18, and holds
at exposed post-recovery states under atomic or recoverable dual publication.

No silent policy substitution is a composite obligation rather than one
standalone theorem. The binding must commit to \(\nu_B\) and its policy inputs,
verification must refer to that same binding, collapse authorization must check
the current or grandfathered policy, and Chapter 16's conflict rule must force
rebinding when the contract changed. No in-place revision of historical
decisions follows abstractly from append-only transition rules; in a physical
implementation, collision-resistant chaining makes unauthorized revision
detectable against an authentic anchor, while actual prevention requires
append-only storage or equivalent replication controls.

\section{The residual sets, recalled}
\[
 \mathcal R_{\mathrm{block}}
 =\{O_i\in\mathcal O\mid a_i=\Required
       \land s_i\neq\mathsf{SAT}\}.
\]
\[
 \mathcal R_{\mathrm{waived}}
 =\{O_i\in\mathcal O\mid a_i=\Waived
       \land s_i\neq\mathsf{SAT}\},
\]
Collapse requires \(\mathcal R_{\mathrm{block}}=\varnothing\) and places no
eligibility condition on \(\mathcal R_{\mathrm{waived}}\). Every waived
deficit nevertheless remains disclosed with its evidential status and
authorization.

\section{Proof obligations and status}
\textbf{Synthesis only.} Typing and guard consequences are unconditional only
within the specified transition system. Ledger completeness, policy continuity,
history uniqueness, and tamper detection require their named publication,
versioning, linearizability, encoding, hash, and authentic-anchor hypotheses.
Canonical encoding, domain separation, fail-closed verifier reporting,
append-only storage, protocol-appropriate honest consensus quorums, and
external actuation contracts remain implementation assumptions. Appendix C
records each claim in the appropriate class.

\chapter{Liveness and Refusal Deadlock}
\section{Chapter thesis}
A system that collapses nothing violates no collapse invariant and remains
useless. Permanent quarantine, impossible repairs, cyclic evidential
dependencies, unavailable verifiers, and policy churn can all preserve safety
while preventing resolution. Liveness therefore requires progress hypotheses
that safety alone cannot supply.

\section{A minimal counterexecution}
Let a required obligation demand evidence that does not and will never exist.
Strict discharge gating correctly prevents collapse forever. The architecture
therefore cannot promise both strict discharge and eventual admission of every
request. The defensible liveness target is eventual terminal resolution:
collapse or final refusal.

\section{Proof obligations and status}
\textbf{Theorem (Conditional Liveness).} Assume: (L1) every binding has a
finite obligation set; (L2) each root permits at most a fixed total number of
ordinary reopening attempts, not merely bounded depth along one branch; (L3)
bounded fairness, so every pending item is scheduled within a known bound; (L4)
every invoked verifier returns a typed verdict within a known bound; (L5) each
policy and binding specification remains stable for a known interval long
enough for one end-to-end attempt; (L6) exhausted ordinary repair triggers an
escalation procedure that terminates within a known bound in admission or
final refusal; and (L7) every non-verifier internal stage completes or returns
a typed failure within a known bound. Then every submitted event reaches
successful collapse or final refusal within a finite computable bound.

\emph{Proof sketch.} By (L2), only finitely many ordinary attempts can occur.
By (L1), (L3), (L4), (L5), and (L7), each attempt is scheduled and completed
within a bounded interval rather than being starved, stuck, or perpetually
invalidated by policy churn. It either collapses or consumes one of the finite
ordinary attempts. Exhaustion invokes escalation, which terminates by (L6).
Summing the per-attempt, scheduling, and escalation bounds yields a finite
global bound. \(\blacksquare\)

Ordinary eventual fairness would justify eventual resolution but not the
stronger phrase “within bounded time”; bounded fairness and bounded stage
execution are what license that conclusion. The result resembles a CAP-shaped
tradeoff only informally. It is not an application of the CAP theorem
\cite{gilbert-lynch}.

\chapter{Threat Model}
\section{Trust boundary}
Sources may submit adaptive malicious inputs. A compromised kernel may diverge
from the artifact named by \(\nu_K\) if deployment attestation fails. A
malicious verifier may fabricate discharge vectors rather than merely share a
correlated bug. Authorities may issue illegitimate capabilities or policies;
clocks may be skewed; storage replicas may suppress, rewrite, or equivocate;
and the network may delay, reorder, duplicate, or drop messages.

\section{Categories of failure}
Byzantine behavior is arbitrary and adaptive \cite{lamport-byzantine}.
Ordinary faults are crashes,
timeouts, partitions, and hardware failures without malicious intent.
Specification error occurs when honest components faithfully implement an
inadequate requirement. Insider abuse uses legitimately held authority for an
improper purpose, and therefore calls for separation of duties and audit rather
than authentication alone. Institutional capture occurs when the policy-making
body itself comes to serve interests contrary to the governed domain. The
architecture cannot cryptographically rescue a captured standard; it will
enforce that standard rigorously.

\section{Proof obligations and status}
\textbf{Load-bearing specification chapter.} Chapter 13 permits Byzantine
sources and network manipulation but excludes compromise of the signing keys
and atomic-head arbiter. Chapter 16 abstracts its trust boundary as
linearizability; a consensus realization must retain whatever honest quorum
its particular protocol requires, not a generic unspecified majority.
Chapter 18 assumes collision resistance and at least one authentic later
anchor outside the adversary's rewrite power. These are hypotheses, not a new
theorem. Their centralization here prevents one chapter's trust boundary from
being silently imported into another.

\chapter{Recovery and Replay}
\section{Chapter thesis}
Chapter 18 proved crash-safe dual publication at collapse. Recovery across the
whole pipeline requires a further distinction: reproducing a deterministic
computation is not the same as re-deciding a historical transition.
Authentication, refusal, authorization, clock reads, and concurrency outcomes
are historically contingent. Recovery must normally reapply their durable
records, not ask the current world to decide them again.

\section{Write-ahead persistence at every boundary}
POP must durably witness successful intake or salvage failure before releasing
an envelope downstream. REFUSE must persist its decision before communicating
or acting upon the disposition. BIND may recompute obligations only when every
historical input, including \(\Gamma_t\), policy, and \(\nu_B\), is durably
available; otherwise it replays the recorded binding. TRANSFORM may rerun under
Chapter 10's replay-equivalence hypotheses, but applying a previously recorded
candidate is sufficient when its witness is already durable.

VERIFY may be recomputed only under an independently stated deterministic
verifier discipline. Otherwise recovery replays the signed verification record
as historical output and never upgrades it by consulting a newer verifier.
COLLAPSE uses Chapter 18's write-ahead transaction. PREPARE, ACTUATE, and
FINALIZE use Chapter 17's pending-intent and attempt journals: a prepared intent
with no attempt witness is pending first actuation, while an attempt without a
conclusive receipt is ambiguous and must enter reconciliation. Multiple
receipts sharing one intent key are records concerning one requested effect;
whether the external system produced one effect or several remains a
reconciliation question, not something key equality alone proves.

\section{Derived indexes}
Forward-lineage indexes and ledger projections are caches over canonical
records. They may be rebuilt by scanning the ledger and need not block recovery
unless an operation depends on them for a safety check. An index used only for
query acceleration can return later; an index used to enforce nonce spending
or uniqueness is safety-critical state and cannot be dismissed as merely
derived without a validated reconstruction first.

\section{Proof obligations and status}
\textbf{Theorem (Crash-Recovery Equivalence).} Let a checkpoint commit to a
canonical machine state and a ledger sequence number. Assume: (R1) every
post-checkpoint transition that became externally visible has a complete,
authenticated durable transition record or a recoverable write-ahead record;
(R2) the durable prefix and record order remain intact; (R3) the state-application
function for each record kind is deterministic; (R4) application is idempotent
under the record's stable identifier; (R5) every historical artifact needed to
validate a record remains available; and (R6) recovery completes before normal
readers and writers observe the reconstructed state. Then replaying the ordered
durable suffix reconstructs the same canonical internal state that those
recorded transitions produced before the crash, excluding explicitly
unreconciled external-world outcomes.

\emph{Proof sketch.} The checkpoint supplies the base state. Induct over the
ordered durable suffix. At each step, (R3) applies the same authenticated record
to the same preceding canonical state and therefore yields the same successor.
If the transition was already applied before the crash, (R4) makes replay a
no-op; if it was only prepared, its write-ahead status determines which
application remains. Hypotheses (R1), (R2), and (R5) prevent missing,
reordered, or unverifiable inputs, while (R6) hides intermediate states in
which only part of a multi-store publication has been restored. Induction
therefore reconstructs the same canonical internal state. An external effect
with no conclusive receipt is not reconstructed by assertion and remains in
reconciliation. \(\blacksquare\)

This theorem intentionally requires deterministic record application, not
deterministic historical decision-making. Rerunning POP would generate a new
observation time; rerunning REFUSE under current policy could revise the past;
rerunning COLLAPSE would enter a new concurrency race. Recovery preserves what
the system durably decided, while selective recomputation is permitted only
for transitions, such as a suitably confined TRANSFORM, whose replay theorem
has separately been established.

\chapter{Observability and Restorability}
\section{Chapter thesis}
Logs, metrics, and traces are not sufficient merely because they exist.
Operational observability is the degree to which the witness plane supports
restoration, explanation, and re-evaluation. Restoration reconstructs canonical
state under Chapter 23's model. Explanation reconstructs why a decision was
made under the policy and evidence then available. Re-evaluation applies a
declared current procedure without rewriting the historical result.

\section{Dependency-closed witness sets}
For a scope \(U\), define \(W(U)\) as the transitive closure of every canonical
record concerning \(U\), every parent and input commitment those records name,
and every versioned policy, binding specification, kernel, verifier, witness,
and checkpoint required by the requested operation. Sufficiency is
purpose-relative. A hash commitment may suffice to check continuity but not to
explain content; a redacted payload may make restoration possible while making
substantive re-evaluation impossible.

\section{Local and global restoration}
A local set closes over one event's repair descendants and their referenced
artifacts. Under Chapter 21's bound on the total number of ordinary attempts
and a bounded escalation protocol, the number of transition records in this
closure is bounded by the protocol constants. Its byte size is not thereby
bounded: payloads, traces, proofs, and artifacts can vary arbitrarily unless
separate size limits are imposed.

Global crash restoration does not require the entire ledger from genesis when
a trusted canonical checkpoint exists. Its sufficient set is the latest valid
checkpoint, the ordered durable suffix after it, and the historical artifacts
needed to validate and apply that suffix. Full historical explanation or audit
may still require the earlier ledger. Recovery sufficiency and archival
completeness are therefore different properties.

\section{Proof obligations and status}
\textbf{Supporting characterization.} Dependency closure is sufficient under
the machine model because every transition reads only its declared inputs and
Chapter 23 recovers by applying authenticated records. Necessity is relative to
the requested purpose: removing a component is necessary exactly when no
remaining record or valid derivation reconstructs the information it supplies.
Prediction of future failures and detection of institutional capture are
research conjectures, not consequences of this characterization.

\part{Implementation and Evaluation}

\chapter{A Reference Data Model}
\section{Chapter role}
The abstract tuples now receive concrete fields and representation invariants.
The following condensed schema preserves the identity and replay corrections
made in Chapters 5, 10, 11, 15, and 19.

\begin{scriptsize}
\begin{verbatim}
EventCore { u, source, content_commit, payload_handle }
Event     { id, u, observed_at, source, content_commit,
            payload_handle, metadata, parent, creation_rank }
Binding   { event_id, event_commit, certificate, obligations, nu_b }
Obligation{ id, predicate, method, witness_class, discharge_policy }
Candidate { output_commit, binding_commit, history_head,
            execution_witness, trace_ref, capabilities, nu_k }
VerificationCore { verdict, candidate_commit, binding_commit,
                   discharge_vector, r_block, r_waived, witness, nu_v }
CollapseProposal { candidate_commit, verification_commit,
                   expected_head, sigma_c, nonce }
CommitCore { id_c, prior_head, candidate_commit, verification_commit,
             sigma_c, nonce, committed_at, authority }
LedgerCore { seq, kind, subject, payload_commit, policy,
             parents, recorded_at, actor, prior_record_hash }
ResolverRecord { payload_handle, disposition, authorization, recorded_at }
Signed<T> { object: T, signer, algorithm, signature }
\end{verbatim}
\end{scriptsize}

The identifier is
\[
 id(e)=\Hash(\mathsf{domain}_e\parallel
             \operatorname{encode}(\langle u,s,h,p\rangle)).
\]
Observation time and metadata remain outside the identity core, while the
payload handle \(p\) is immutable. Redaction changes a \(ResolverRecord\), not
the event. Candidate records retain both \(h(b)\) and the history head they
read; verification records retain \(h(x)\) and \(h(b)\), making
\(\operatorname{SameCandidate}\) representable.

\section{Canonical encoding}
Fields are encoded in fixed schema order. Every variable-length component is
length-prefixed, every sum type carries an explicit variant tag, maps and sets
are sorted by declared canonical keys, scalar encodings are unique, and Unicode
and numeric normalization rules are versioned. These rules together, not field
order alone, make encoding injective over well-typed normalized values.

Each object kind has a fixed reserved domain tag. Appendix D is the registry.
An identifier is computed over an unsigned core; authentication then wraps
that core in \(\operatorname{Signed}\langle T\rangle\). A signature is never
included in the bytes whose identifier it authenticates, avoiding circular
self-reference. Ledger signatures likewise authenticate \(LedgerCore\), and
commit signatures authenticate \(CommitCore\).

\section{Migration and payload storage}
Schema migration produces a versioned derived view and never rewrites original
bytes. Payload handles may be ordinary opaque locators, randomized
content-addressed capabilities, or keyed addresses. A public locator of the
form \(h(payload)\) is unsuitable for low-entropy secret payloads because it
enables dictionary testing; integrity is checked against the event's commitment
after authorized retrieval instead.

\section{Proof role}
\textbf{Implementation-bearing.} The schema exhibits a realizable witness for
Chapter 5's encoding and domain-separation hypotheses, conditional on the leaf
encoders and normalization rules being injective. Cross-language golden vectors
must confirm that independent implementations produce identical bytes and
identifiers.

\chapter{A Minimal Spherepop Service}
\section{Pure decisions and effectful adapters}
The service separates transition decisions from persistence and external
effects. A pure transition receives explicit state, authenticated facts, clock
values, and inputs, and returns a typed result plus records to publish. It does
not read a live clock, database, network, or random source. POP's authentication
result and observation time are therefore explicit inputs; COLLAPSE can
construct and validate a proposed update purely, while only the persistence
adapter can perform the linearizable compare-and-append.

\begin{scriptsize}
\begin{verbatim}
bind(event, admission, constraints,
     environment, binding_version)
  -> Bound(binding)
   | Deferred(reason)
   | Conflict(reason)
\end{verbatim}
\end{scriptsize}

The adapter implements durable witnesses, write-ahead publication, nonce and
attempt journals, checkpointing, and the external PREPARE--ACTUATE--FINALIZE
protocol. Keeping decision logic separate makes failures of storage semantics
distinguishable from failures of transition logic.

\section{Service boundaries and idempotency}
Submission should be asynchronous by default: intake returns an identifier,
while later status or notification reports refusal, repair, escalation, or
commitment. A synchronous facade may wait for fast terminal cases but cannot
make instant resolution part of the protocol contract.

A client idempotency key is mapped durably to the source-scoped occurrence
token \(u\) before envelope construction. It complements event identity without
changing it. Because observation time and metadata are already excluded from
the identity core, client retries do not need to manufacture new identities
merely because transport metadata changed.

Transaction scopes follow safety boundaries. Early transitions atomically
publish their outputs and ledger witnesses within their relevant stores.
COLLAPSE requires the linearizable head operation of Chapter 16 and the
recoverable dual publication of Chapter 18. No claim that an operation is
“local” excuses it from the witness-plane append.

\section{Deterministic fixtures}
Golden fixtures check byte-identical BIND outputs for fixed
\((e,\alpha,C,\Gamma_t,\nu_B)\), and replay fixtures check TRANSFORM under all
of Chapter 10's hypotheses, including pinned state and captured nondeterminism.
Fixtures test the implementation of those hypotheses; they do not replace the
proofs.

\section{Proof role}
\textbf{Refinement obligation.} A mapping \(\alpha\) from implementation states
to abstract states must make every visible service operation simulate a valid
abstract transition, with stuttering steps permitted for internal persistence
work. Full refinement requires models of the actual database isolation level,
message delivery, compare-and-append primitive, recovery protocol, clocks, and
external effect boundary. Documentation and tests provide conditional evidence;
they are not a completed simulation proof.

\chapter{Testing the Boundaries}
\section{Seven test families}
Property-based tests exercise BIND determinism, verdict partition and error
normalization, waiver independence, and encoding round trips over generated
well-typed inputs. State-machine tests generate transition sequences and check
ledger correspondence, sole-writer behavior, lineage acyclicity, and nonce
spending after every step.

Bounded model checking exhaustively explores small arbiter and recovery models,
especially competing collapses against one head. Its completeness applies only
to the stated finite model and bound. Fuzzing attacks POP parsing, canonical
decoding, signature envelopes, and typed failure paths with malformed bytes.

Fault injection crashes the service before and after each durable boundary,
testing Chapter 23's complete-record premise and idempotent application
requirement. Concurrency tests exercise the deployed sequencer or consensus
mechanism rather than its simplified model. Adversarial fixtures instantiate
certificate replay, substitution, stale authorization, rollback, resolver
tampering, forged waiver, and equivocation attempts.

\section{Proof role}
\textbf{Empirical support, not proof.} Passing samples does not establish a
universal theorem, and bounded exploration does not establish unbounded
correctness. The value of the mapping is diagnostic: a failure names the
specific invariant, assumption, transition boundary, and durable record whose
implementation diverged from the abstract machine.

\chapter{Evaluation Programme}
\section{Metrics}
Measure throughput and latency separately for each stage; audit amplification
as records and bytes in \(L\) relative to commits and bytes in \(H\), treating
empty denominators explicitly; recovery time against checkpoint age and
unapplied suffix length; verification cost by verifier and quorum; stale-head
and commit-conflict rates; repair success; indeterminate and error verdict
rates; escalation frequency; and storage, replay, and validation cost for each
execution-witness style.

The proposed “false refusal” metric is rejected because it contradicts Chapter
8: a descendant's later success does not make the parent's earlier refusal
false. The measurable quantity is \emph{first-pass repair burden}: the fraction
of ultimately successful lineages that required at least one refusal and
repair, accompanied by an audit of whether each original refusal was correct
under its contemporaneous policy. Burden measures friction without rewriting
historical truth.

High audit amplification may reflect either pathological churn or appropriately
rich evidence; high stale-head frequency may motivate sharding; low repair
success may reveal decorative repair paths; frequent indeterminacy may reveal
unobtainable witness classes. Metrics diagnose questions rather than answer
them by themselves.

\section{Methodological discipline}
Confirmatory comparisons preregister workloads, metrics, baselines, exclusion
rules, and fault injections before results are inspected. Continuous operational
monitoring need not be preregistered, but any later inferential claim must
distinguish exploratory discovery from confirmatory evaluation.

\section{Proof role}
\textbf{Empirical chapter.} Performance and assurance are orthogonal. A fast
non-linearizable arbiter remains unsafe, and a low repair burden does not prove
that refusal policy is just. Benchmarks describe observed costs and frequencies;
they establish none of the formal properties in Chapters 4 through 24.

\chapter{Three Applied Case Studies}
\section{Scientific claims under unsettled evidence}
A manuscript and its supporting materials form an event. Journal intake is
POP; editorial disposition is REFUSE's total gate, returning either admission
or a typed refusal. Revision is REOPEN: a new submission with a new occurrence
token and explicit lineage. BIND fixes methodological, disclosure, and
statistical obligations before the reported candidate is assessed. Registered
protocols, versioned datasets, executable analysis, and environments supply a
domain analogue of \(\nu_K\). Peer and statistical review instantiate VERIFY,
with \(\Indeterminate\) expressing that the available evidence does not settle
a required question. Publication is COLLAPSE only relative to the journal's
own authoritative publication history, not to the field's globally accepted
beliefs.

Later replication failure may justify correction, retraction, or a new study.
It does not erase what reviewers saw or which standards governed the earlier
decision, but neither does historical preservation certify that the earlier
review was correct relative to its own evidence. Fraud, oversight, or
misapplication may have made it defective even then.

The mapping imports fallible human review, domain standards, and a journal-local
arbiter. Science as a whole has no linearizable global head, so History Never
Forks does not apply to the literature. Human reviewers also do not reliably
self-report their own failure modes in the manner Chapter 11's trusted wrapper
requires.

\section{A distributed software transformation}
A pull request is an event; continuous-integration intake, policy gates,
required tests, reproducible builds, test verdicts, and merge correspond to
POP, REFUSE, BIND, TRANSFORM, VERIFY, and COLLAPSE. Build identity pins the
container, compiler, configuration, dependency lockfile, and relevant
environment. A flaky test belongs under \(\Indeterminate\) or
\(\ErrorVerdict\), depending on whether the test completed inconclusively or
the checking machinery failed. A logged emergency override is a waiver, not a
silent conversion of UNSAT into SAT. An outdated pull request is the familiar
stale-head case.

Version-control hashes are close analogues of domain-separated content
commitments, but they are not literally Chapter 5's event identifiers: commit
objects commonly include parents, authorship metadata, and messages, whereas
Spherepop identity includes a source-scoped occurrence token and immutable
payload handle. Deployment then invokes PREPARE, ACTUATE, reconciliation, and
FINALIZE rather than treating merge as proof that production changed.

The mapping imports uncompromised build infrastructure and collision-resistant
hashing. The practical SHA-1 collision demonstrated in 2017 shows why hash
agility is an operational requirement rather than a decorative caveat
\cite{stevens-sha1}. Supply-chain compromise violates the kernel or verifier identity
boundary; the architecture names that failure without solving it.

\section{An irreversible actuator}
A command to position an industrial valve is admitted and authenticated at
POP, checked against interlocks and operator capability at REFUSE, and bound to
pressure, conflict, authorization, and sensor obligations. TRANSFORM simulates
a pressure and flow trajectory without touching the valve. VERIFY checks that
candidate, returning \(\Indeterminate\) when uncertainty straddles a safety
threshold.

A successful collapse proposal enters Chapter 17's external-effect protocol:
PREPARE durably records intent, ACTUATE sends the physical command, and
FINALIZE appends the operational commit only after a conclusive receipt or
reconciliation. Flow that already occurred cannot be rolled back; closing the
valve is a new compensating event.

The mapping imports trustworthy sensors, adequate physical models, actuator
identity, and an external duplicate-handling contract. A correct simulation
can omit a decisive physical variable, exactly as Chapter 14 and Interlude III
warn.

\section{Proof role}
\textbf{Model validation, not theorem.} Each case demonstrates a typed
correspondence and exposes its imported assumptions. The formal guarantees
hold only after a domain has genuinely supplied those assumptions; an elegant
mapping cannot establish honest reviewers, uncompromised infrastructure,
adequate models, or trustworthy sensors.

\part{Integration and Consequence}

\chapter{From RSVP to Event Histories}
\section{Standing caution}
RSVP is treated here as a possible continuous scalar-vector substrate and KES
as a provisional account of event-mediated historical fixation. No detailed
KES state vector, admissibility field, or monotonic-contraction law is treated
as canonical \cite{flyxion-rsvp-kes}. In particular, global monotonic contraction is too strong if it
forbids repair and reopening.

\section{Bridge principle}
Spherepop is a candidate worked instance of the abstract pattern KES gestures
toward. RSVP would supply continuous dynamics; KES would name the emergence of
historically fixing events; Spherepop supplies typed computational boundaries
for encounter, admission or refusal, binding, candidate production,
verification, and commitment, with the ledger witnessing each boundary. This
is a correspondence of roles, not a derivation of one theory from another.

The apparent conflict with monotonic contraction can be localized. A refused
event never regains admissibility in place. REOPEN creates a descendant event
with a fresh identity and its own admissibility judgment. Contraction may
therefore be coherent per event-attempt while the global possibility structure
branches through new descendants. The growing audit history \(L\) and selective
operational history \(H\) provide distinct indices that the earlier sketch
lacked.

Any genuine RSVP bridge would still need an observation or discretization map
from continuous field state to an event envelope, and an admissibility map from
field-relative conditions to REFUSE and BIND outputs. Their domains,
codomains, covariance properties, and dimensional consistency remain
unspecified.

\section{Proof role}
\textbf{Speculative integration.} The chapter asserts only that the types admit
a plausible correspondence. RSVP equations do not entail Spherepop
transitions, and Spherepop's theorems do not validate KES mathematics.

\chapter{Distinction Holonomy}
\section{Paths and terminal equivalence}
The terminology and general framing are inherited from the author's separate,
unpublished treatment of distinction holonomy \cite{flyxion-holonomy}; the
reduced definitions here do not validate that larger formalism.
Because a repair graph can branch, a path \(\pi\) is one causally continuous,
ordered root-to-terminal branch together with the transition records generated
along that branch. The entire descendant graph is not one path. Two paths are
terminally equivalent under a declared projection \(T\), written
\(\pi_1\sim_T\pi_2\), when \(T(\pi_1)=T(\pi_2)\). The projection may compare
only success, committed candidate content, or a finer terminal structure.

\section{Transported residue}
Rather than an undefined set-theoretic “complement,” define a declared residue
map \(R_T\) that extracts the path information intentionally ignored by
\(T\): refusals, repairs, elapsed time, verifier identities, policy versions,
waivers, authorities, and other provenance. Terminally equivalent paths may
therefore carry different \(R_T\) values.

An actual holonomy theory would additionally require a transported object, a
transport operator along composable path segments, a notion of closed loop or
alternative-path comparison, and an invariant showing how transport around
that loop fails to return the object unchanged. A distance between residues is
not by itself holonomy, and dependence on an arbitrary projection must be
analyzed rather than hidden.

\section{Proof role}
\textbf{Prospective formalization.} The chapter supplies disciplined
definitions and a list of missing structures. It claims no holonomy theorem.

\chapter{MEM|8 and Resonant Recall}
\section{Archive, salience, and admissibility}
The resonant-recall proposal derives from separate MEM|8 manuscripts
\cite{flyxion-wave-machine,flyxion-conscious-gaia}. It is treated here as a
candidate integration model rather than as a result of Spherepop.
Archive is the canonical \(L\) and \(H\). Salience is a mechanism for ranking
what may be useful to recall. Admissibility determines what a requester and
the retrieval computation itself may access. These are three distinct
relations. A highly resonant record may still be forbidden, and a complete
archive does not imply complete approximate recall.

\section{Explicit retrieval map}
Any derived memory system must declare
\[
 \rho(q,c)=\operatorname{Rank}_{q,c}
             (\operatorname{Eligible}_{c}(L)).
\]
Eligibility must constrain the index or retrieval computation before
unauthorized records are scored. Filtering only the displayed results after a
global ranker has processed secret material is insufficient because scores,
embeddings, caches, and associations can themselves leak information.
Deployments therefore need policy-partitioned indexes, admissibility-preserving
representations, or cryptographic computation that enforces the same boundary.

Ledger completeness and tamper evidence remain properties of the upstream
archive; they do not “survive” as completeness properties of a lossy latent
memory. A resonant layer may omit or distort retrieval and must never be the
only route to canonical records. Exact addressed lookup and verification
remain available through the archive.

\section{Proof role}
\textbf{Speculative integration.} MEM|8 is one candidate salience model.
Spherepop neither entails coupled-oscillator recall nor validates a latent
trajectory formalism. The established contribution is the integration
constraint separating archive, admissibility, and salience.

\chapter{Institutions as Event Systems}
\section{Legal procedure}
Filing resembles POP; threshold dismissal can resemble REFUSE; amendment is a
new descendant rather than mutation; legal elements resemble bound
obligations; and judgment resembles commitment. The mapping must not be made
too exact. Standing is principally an authority or capability question,
jurisdiction concerns the forum's authority and governing policy, and neither
is simply provenance.

A hung jury is not automatically Chapter 11's \(\Indeterminate\). It may
reflect evidential uncertainty, disagreement over interpretation, or a
non-unanimous aggregation rule. An acquittal likewise means the prosecution
did not satisfy its burden, not necessarily that innocence was positively
proved. Legal verdicts need their own typed aggregation semantics before the
four verifier verdicts can be imported.

Vacatur is also subtler than in-place deletion. A well-kept docket preserves
the original judgment and appends the appellate act that removes its
authoritative legal effect. Doctrinally the judgment is vacated; historically
the issuance and vacatur both remain. This resembles descendant correction and
status supersession more closely than erasure, although legal consequences and
precedential treatment exceed Spherepop's present model.

\section{Governance and interpretation}
Bill introduction, committee disposition, procedural constraints, voting, and
enactment offer a rough mapping. Real legislatures, however, renegotiate both
the candidate and the standards for accepting it throughout deliberation.
Consensus bodies may revise the target and proposal together. This is
continuous rebinding, or something beyond the pipeline, rather than verification
against a contract fixed in advance.

The deepest disanalogy is interpretation. Courts and institutions often decide
what “reasonable,” “material,” or “in the public interest” means through the
act of applying it. Spherepop assumes that enough interpretive labor has
already occurred for BIND to emit checkable obligations. It has no theorem
reducing institutional judgment to predicate evaluation.

\section{Proof role}
\textbf{Analogical and diagnostic.} The vocabulary can locate conflated
boundaries, but it neither proves institutional adequacy nor supplies a theory
of interpretation.

\chapter{The Ethics of Preserved Refusal}
\section{The ethical tension}
Permanent preservation supports accountability while creating durable risks
for subjects of unsubstantiated allegations, third parties named incidentally,
and people whose individually correct denials are later aggregated into an
improper adverse profile. Structural metadata and commitments can themselves
be sensitive even after payload erasure, especially when values are
low-entropy or linkable.

\section{Four principles}
Proportional retention assigns different availability and disclosure periods
to unsubstantiated attempts, denied claims, and committed findings without
rewriting the structural ledger. Contestability requires notice,
comprehensible reasons, a usable reopening process, and often practical support;
an API endpoint alone is not meaningful recourse.

Redaction may protect a vulnerable third party or conceal misconduct by a
powerful actor. An authorization record proves that a decision was recorded,
not that it was ethically justified. Access control should therefore default
to directly involved parties and independently accountable audit functions,
with wider disclosure requiring a specific recorded basis. Chapter 18's
projections make these policies implementable but cannot choose among them.

\section{Proof role}
\textbf{Normative chapter.} Completeness, tamper evidence, redaction continuity,
and authorization records do not determine retention duration, default access,
or legitimate purpose. Formal consistency cannot choose the ethical premise.

\chapter{Conclusion: Commitment Without Conflation}
\section{Four promises}
A trustworthy event architecture does not promise that every committed state
is true or good. It makes four narrower promises under named conditions.
Encounter does not masquerade as admission: POP records occurrence and
REFUSE's total gate separately returns admission or a reason-bearing
non-admission. Execution does not masquerade as verification: TRANSFORM
produces a candidate, while VERIFY alone evaluates the obligations fixed by
BIND.

Verification does not masquerade as truth: PASS means that the declared
required obligations were reported SAT under their declared methods and valid
checking relation. It says nothing by itself about the adequacy of the
obligation set or the goodness of the governing policy. Commitment does not
occur without reconstructible warrant: the collapse guard, sole-writer rule,
and witness plane connect every visible post-recovery commit to exactly one
collapse witness under the publication hypotheses.

\section{Open boundary}
Institutional capture, interpretive judgment, ethical retention, RSVP--KES
integration, distinction holonomy, and MEM|8's resonant dynamics remain open.
They are not omissions disguised as future certainty. They mark the boundary
between the architecture's proved properties and the surrounding questions
that require other theories, empirical evidence, or normative argument.

\section{Closing}
The witness plane cannot make policy legitimate or preserved allegations
harmless. It can prevent an attempt from being silently rewritten as success,
a computation from being mistaken for authorization, a passing check from
being inflated into truth, and a commitment from being detached from the
historical warrant that produced it. This is not certainty. It is accountable
commitment, falsifiable at named boundaries and offered as nothing more.

\section{Proof role}
\textbf{Synthesis only.} No new theorem is introduced.

\appendix

\chapter{Notation and Type Catalogue}
\section{State and event types}
\[
 \mathcal S_t=\langle Q_t,L_t,H_t,P_t,A_t,W_t\rangle ,
\]
where \(Q_t\) is pending work, \(L_t\) the audit ledger, \(H_t\) operational
history, \(P_t\) the policy and version registry, \(A_t\) authority state, and
\(W_t\) durable protocol state containing write-ahead transactions, pending
external intents, actuation attempts, receipts, checkpoints, and spent nonces.

An event has immutable identity core and full envelope
\[
 c(e)=\langle u,s,h,p\rangle,qquad
 e=\langle id,u,t,s,h,p,m,parent,rank\rangle ,
\]
\[
 id(e)=\Hash(\mathsf{domain}_e\parallel\operatorname{encode}(c(e))).
\]
Here \(u\) is the source-scoped occurrence token, \(t\) system observation
time, \(s\) authenticated source, \(h\) randomized or otherwise
policy-appropriate content commitment, \(p\) immutable payload handle, and
\(m\) non-causal metadata. Redaction changes
\(\mathsf{Resolve}_t(p)\), never \(p\) or the envelope.

\section{Transitions}
\[
\begin{aligned}
 \Pop(e_{\mathrm{raw}}) &\to \mathsf{POP\_RESULT},\\
 \Refuse_{\sigma_t}(e,H_t) &\to
   \mathsf{ADMIT}(\alpha)\mid\mathsf{REFUSAL}(d),\\
 \operatorname{REOPEN}(id(e),\Delta',w,u') &\to e',\\
 \Bind(e,\alpha,C,\Gamma_t,\nu_B) &\to
   \mathsf{BOUND}(b)\mid\mathsf{DEFERRED}(d_B)\mid\mathsf{CONFLICT}(f_B),\\
 \Transform_{K,\theta}(b,H_t) &\to
   \mathsf{CANDIDATE}(x)\mid\mathsf{FAILURE}(f_K),\\
 \Verify_{V,\sigma_v}(x,b,H_t) &\to v,\\
 \Collapse(c,H_t,L_t,W_t) &\to \mathsf{COLLAPSE\_RESULT},\\
 \operatorname{PREPARE}(c)&\to\iota,\\
 \operatorname{ACTUATE}(\iota)&\to\epsilon\mid\mathsf{AMBIGUOUS},\\
 \operatorname{FINALIZE}(\iota,\epsilon)&\to\kappa .
\end{aligned}
\]
Every branch writes a typed audit witness. Early-stage transitions do not
write \(H\). External-effect finalization invokes the same successful commit
publication primitive as direct collapse; it is not a second independent
writer.

\section{Refusal, obligations, and verdicts}
\[
 d=\langle r,\rho,\Delta,\tau,\pi,terminal\rangle ,
\]
with \(r\) in quarantine, repair, evidence-only, defer, redact, or deny.
Terminality is an orthogonal Boolean set only by the bounded escalation
procedure; it is not a seventh disposition. Reason grounds are identity,
schema, provenance, policy, capability, and temporal.

\[
 O_i=\langle P_i,m_i,w_i,\delta_i\rangle,\quad
 d_i=\langle s_i,a_i\rangle ,
\]
where \(s_i\) is \(\mathsf{SAT}\), \(\mathsf{UNSAT}\), \(\mathsf{UNKNOWN}\), or
\(\mathsf{NOT\_CHECKED}\), and \(a_i\) is \(\mathsf{Required}\)
or \(\mathsf{Waived}\). The blocking and disclosure residuals are
\[
 \mathcal R_{block}=\{O_i:a_i=\Required\land s_i\neq\mathsf{SAT}\},
 \quad
 \mathcal R_{waived}=\{O_i:a_i=\Waived\land s_i\neq\mathsf{SAT}\}.
\]
The overall verdict \(q\) is PASS, FAIL, INDETERMINATE, or ERROR. ERROR
normalizes every \(s_i\) to \(\mathsf{NOT\_CHECKED}\).

\section{Candidate, verification, and collapse}
A candidate commits to output, binding, and read head:
\[
 x=\langle y,h(b),h(H_t),\chi,T,\mathcal C,\nu_K\rangle .
\]
A verification core contains \(q, h(x), h(b),\mathbf d,\zeta,
\mathcal R_{block},\mathcal R_{waived},\nu_V\) and is authenticated by a
separate signature envelope. Thus
\[
 \operatorname{SameCandidate}(x,v)
 \iff h(x)=v.h_x\land h(b)=v.h_b.
\]

\[
 c=\langle h(x),h(v),h(H_t),\sigma_c,\eta\rangle ,
\]
and the unsigned commit core is
\[
 \kappa=\langle id_c,h(H_t),h(x),h(v),\sigma_c,\eta,t_c,a_c,
                 \mathcal R_{waived}\rangle .
\]
Its signature envelope is external to the core, and the resulting head is
\[
 h(H_{t+1})=\Hash(\mathsf{domain}_H\parallel h(H_t)\parallel
                         \operatorname{encode}(\kappa)).
\]
Collapsibility requires PASS, empty \(\mathcal R_{block}\), freshness, and
candidate identity. Freshness has elapsed-time, state-head, policy, capability,
and evidence components.

\section{Hypothesis catalogue}
Encoding hypotheses are injective canonical encoding and exclusive domain
separation. History hypotheses are unique genesis and linearizable
compare-and-append. Liveness hypotheses \((L1)\)--\((L7)\) are finite
obligations, a bound on total ordinary attempts per root, bounded fairness,
bounded verifier termination, a sufficiently stable policy interval,
terminating escalation, and bounded completion or typed failure for every
other internal stage.

Recovery hypotheses \((R1)\)--\((R6)\) are complete durable records for every
visible transition, durable ordered-prefix integrity, deterministic record
application, idempotent application by stable record identifier, availability
of historical validation artifacts, and recovery isolation before state is
exposed. These replace the earlier and stronger assumption that historical
decisions themselves can always be recomputed deterministically.

\section{Speculative notation}
A distinction path \(\pi\) is a single causally continuous root-to-terminal
branch. Terminal equivalence is \(\pi_1\sim_T\pi_2\) iff
\(T(\pi_1)=T(\pi_2)\); \(R_T\) is the declared residue projection. No
criterion-independent transported complement or holonomy operator is assumed.

\chapter{Operational Semantics}
\section{Frame and publication conventions}
Any state component absent from a rule's conclusion is carried unchanged.
Every transition first creates an authenticated durable transition record,
then applies that record idempotently to canonical state. The notation
\(\mathsf{emit}(r)\) appends the audit witness for record \(r\). The operation
\(\mathsf{commit}(w_\kappa)\) is the sole abstract writer of \(H\): direct
COLLAPSE and external FINALIZE both call this one primitive.

\section{Early stages}
Representative success rules are
\[
\frac{\langle\Pop,e_{raw}\rangle\in Q_t\quad
      \operatorname{normalize}(e_{raw})=e}
     {Q_{t+1}=Q_t-\{\langle\Pop,e_{raw}\rangle\}
        \cup\{\langle\Refuse,e\rangle\},\quad
      L_{t+1}=\mathsf{emit}(\ell^{pop})}
 \;[\mathrm{POP}]
\]
and
\[
\frac{\langle\Refuse,e\rangle\in Q_t\quad
      \Refuse_{\sigma_t}(e,H_t)=\mathsf{ADMIT}(\alpha)}
     {Q_{t+1}=Q_t-\{\langle\Refuse,e\rangle\}
        \cup\{\langle\Bind,e,\alpha\rangle\},\quad
      L_{t+1}=\mathsf{emit}(\ell^{admit})}
 \;[\mathrm{ADMIT}].
\]
The refusal branch removes the work item, appends \(\ell^{ref}\), and schedules
no downstream computation. POP failure and duplicate arrival likewise append
typed records; a duplicate arrival does not create a second event node.

BIND emits Bound, Deferred, or Conflict. Only Bound schedules TRANSFORM.
TRANSFORM emits Candidate or Failure; only Candidate schedules VERIFY. VERIFY
always emits a full typed result, and only PASS with an empty blocking residual
may schedule a collapse attempt. By the frame convention all these rules
satisfy \(H_{t+1}=H_t\).

\section{Commit publication}
For an internal effect, a successful guard creates a durable transaction
\[
 w_\kappa=\langle\kappa,\ell_\kappa^{collapse}\rangle
\]
and invokes \(\mathsf{commit}(w_\kappa)\), which atomically, or through the
recoverable protocol of Chapter 18, appends one \(\kappa\) to \(H\) and one
witness to \(L\). A failed guard appends only a typed collapse-failure record.

For an external effect, PREPARE places \(\iota\) in durable protocol state and
appends its witness. Before transmission, ACTUATE appends an attempt record.
A conclusive receipt is stored; a missing receipt marks the intent ambiguous.
FINALIZE validates \((\iota,\epsilon)\), constructs \(w_\kappa\), and calls the
same \(\mathsf{commit}\) primitive. Consequently the sole-writer premise is
preserved even though two routes reach it.

\section{Recovery rule}
Given checkpoint \(C_j\) and authenticated ordered durable suffix
\((r_{j+1},\ldots,r_n)\), recovery folds the deterministic application
function:
\[
 \mathcal S_n=
 \operatorname{foldl}(\mathsf{apply},C_j,
                     [r_{j+1},\ldots,r_n]).
\]
Already-applied identifiers are no-ops. Historically contingent decisions are
not rerun. A prepared or attempted external intent without a conclusive receipt
enters reconciliation rather than being guessed into success or failure.

\chapter{Core Invariants and Proof Obligations}
This appendix is the canonical registry of formal claims. Each entry must give
its complete hypotheses, conclusion, proof, dependencies, failure conditions,
and prospective mechanization status. It must not flatten the architecture's
claims into a uniform sequence of nominal theorems.

\section{Class I: Provable as stated in the transition system}

\subsection{Early-Stage Non-Interference}
For POP, REFUSE, BIND, TRANSFORM, and VERIFY, prove by inspection of the
small-step rules that no transition writes operational history:
\[
  H_{t+1}=H_t.
\]
The result is load-bearing despite its short proof because every later safety
argument depends on COLLAPSE being the sole writer of \(H\).

\subsection{Layered Admission Non-Implication}
Provide the complete counterexample matrix showing that epistemic,
computational, operational, and irreversible admission are pairwise
non-implying.

\subsection{Identifier Collision Reduction}
Reduce an event-identifier collision to either non-injective canonical encoding
or a collision in \(\Hash\). Do not characterize this reduction as a new proof
of hash collision resistance.

\subsection{Obligation-Generation Well-Definedness}
Prove deterministic obligation generation for well-typed
\((e,C,\Gamma_t,\nu_B)\). Limit the claim to syntactic checks and declared
decidable fragments; general semantic satisfiability remains outside the
result.

\subsection{Verdict Partition and Error Containment}
Prove that the four verification verdicts are exhaustive and exclusive under
the verifier operational semantics. Prove separately that
\(q=\ErrorVerdict\) prevents every \(s_i\) from retaining
\(\mathsf{SAT}\).

\subsection{Collapse Guard Respects Waivers}
Prove that only \(\mathcal R_{\mathrm{block}}\) controls collapse eligibility,
while every member of \(\mathcal R_{\mathrm{waived}}\) remains disclosed with
its evidential status and waiver authority.

\subsection{Ledger Completeness}
Prove by induction that each \(\kappa\in H_t\) has exactly one corresponding
collapse witness in \(L_t\), at every exposed quiescent state, using the
single \(\mathsf{commit}\) primitive and atomic or recoverable dual
publication. Intermediate recovery states need not satisfy the pointwise
invariant and must not be externally exposed.

\subsection{Hash-Chain Tamper Evidence}
Given an authentic trusted later head and collision-resistant hashing, prove
that an altered record cannot validate to that fixed anchor except with
negligible collision probability. If the suffix is honestly recomputed, every
descendant hash changes; detection still requires the segment or an
authenticated proof path.

\subsection{Redaction Preservation}
Prove that authorized payload unavailability leaves \(id(e)\) and \(h(e)\)
unchanged when the canonical envelope, immutable handle \(p\), and content
commitment are retained. Redaction changes resolver state, not \(p\).
Record non-repudiation of erased plaintext as a non-preserved property.

\section{Class II: Provable only under named additional hypotheses}

\subsection{History Never Forks}
Assume a unique genesis history and a linearizable compare-and-append primitive
implemented by a single arbiter or equivalent consensus mechanism. Prove by
induction that at most one candidate can extend a given current head and that
the operational history remains unique.

\subsection{Transformation Replay Equivalence}
Assume deterministic kernels, canonical inputs, pinned kernel and environment
identity, equal initial state, and no undeclared nondeterminism. Prove equality
of replayed outputs up to the declared trace equivalence.

\subsection{Crash-Recovery Equivalence}
Assume complete authenticated records for visible transitions, an intact
ordered durable prefix, deterministic and idempotent record application,
available historical validation artifacts, and recovery isolation. Prove that
checkpoint plus suffix reconstructs the same canonical internal state,
excluding unreconciled external effects. Historical decisions are not rerun.

\subsection{Freshness and Substitution Resistance}
Under collision resistance, signature unforgeability, nonce uniqueness,
correct freshness predicates, and atomic current-head comparison, bound the
probability that a probabilistic polynomial-time adversary commits a stale or
substituted candidate.

\subsection{Conditional Liveness}
Under finite obligations, a bound on total ordinary attempts per root, bounded
fair scheduling, bounded verifier termination, sufficiently stable policy
intervals, bounded completion or typed failure for other stages, and
terminating escalation, prove terminal resolution within a computable bound.
Ordinary eventual fairness supports eventuality, not a numerical bound.

\subsection{ACTUATE Idempotency}
Treat external idempotency as a capability-contract assumption. Spherepop can
prove stable intent, pre-transmission attempt journaling, and idempotency-key
reuse, but it cannot derive the
external system's behavior from its internal transition rules.

\section{Class III: Negative results, limitations, and non-claims}

\subsection{Correlated Verifiers}
Show that \(k\) verifiers sharing a dependency gain no protection against a
fault in that dependency. Do not multiply marginal error rates without an
independence argument.

\subsection{Verification Is Not Truth}
VERIFY establishes obligation satisfaction relative to a declared policy and
verifier. It does not certify that the obligation set is complete, adequate,
or ethically sound. Present counterexamples rather than an irrelevant appeal
to G\"odel.

\subsection{Tamper Evidence Is Not Correctness}
Hash chaining does not prove that an original record was truthful, that all
observers received the same head, that data remain available, or that recorded
times are honest.

\subsection{Compensation Is Not Reversal}
An irreversible external effect may be answered by a later compensating act,
but the original effect is not thereby removed from history.

\section{Mechanization plan}
Formalize the abstract transition system first, then Class I invariants. Add
the linearizable-CAS model before attempting History Never Forks. Encode
cryptographic results as reductions to assumptions rather than reimplementing
cryptographic proofs. Maintain a machine-readable dependency table connecting
each theorem to its hypotheses and the chapters that cite it.

\chapter{Canonical Serialization and Hashing}
\section{Encoding requirements}
For each schema version, \(\operatorname{encode}\) is pure and injective over
well-typed normalized values. Field order and scalar representations are fixed;
sum variants and object versions carry explicit tags; variable-length fields
use fixed-width unsigned length prefixes; maps and sets are canonically sorted;
Unicode, timestamps, integers, and floating-point values have one admitted
normal form. A four-byte length prefix is sufficient only when the schema also
rejects values beyond its representable bound.

\section{Domain registry}
The initial registry assigns distinct tags to event core, refusal core,
admission certificate, reopening record, obligation, binding, candidate,
verification core, collapse proposal, commit core, ledger core, external
intent, actuation attempt, receipt, resolver record, checkpoint, and
write-ahead transaction. Concrete hexadecimal values are controlled by the
registry artifact, not by prose examples. A semantically incompatible schema
version receives a new versioned tag; tags are never silently reassigned.

\section{Hash and signature agility}
\[
 \mathsf{HashVal}=\langle algorithm,digest\rangle .
\]
Equality requires equal algorithm identifiers and digest bytes. A migration
record commits under the old chain to the first head under the new algorithm,
and the new chain commits back to the old anchored head. Without this explicit
bridge, tamper evidence holds only within each algorithm segment.

\[
 \begin{aligned}
 \operatorname{Sign}(sk,T)=\operatorname{Sig}\bigl(sk,{}&
 \mathsf{context}\parallel\mathsf{domain}(T)\\
 &\parallel\mathsf{version}(T)\parallel\operatorname{encode}(T)\bigr).
 \end{aligned}
\]
The signature wraps an unsigned core. Neither the signature nor an identifier
derived from that core is recursively included in its own preimage.

\section{Normative test vector shape}
An event identity vector encodes fields in the order \((u,s,h,p)\), each
preceded by its length and governed by its leaf-type encoding. Changing
\(t\) or \(m\) leaves \(id(e)\) unchanged but changes the full envelope hash.
Changing \(p\) changes event identity; redaction tests instead change
\(\mathsf{Resolve}_t(p)\) and assert that both \(p\) and \(id(e)\) remain
unchanged. Cross-language fixtures must cover empty fields, maximum lengths,
Unicode normalization, map order, every sum variant, unknown-version rejection,
and each algorithm-migration boundary.

\chapter{Schemas and Example Records}
\section{Duplicate-charge lineage}
The normative example follows Interlude I. Values below are schematic;
production encodings use full digests, immutable occurrence tokens, and signed
cores.

\begin{scriptsize}
\begin{verbatim}
POP {
  event: { id:"e0", u:"u0", source:"agent:priya",
           content_commit:"h0", payload_handle:"p0",
           observed_at:"2026-09-01T14:02:07Z" },
  dedup:false
}
REFUSAL {
  subject:"e0", disposition:"REPAIR", reason:"schema",
  repair:{ required:"original_transaction_id" },
  evidence_commit:"pi0", terminal:false
}
REOPEN {
  event:{ id:"e1", u:"u1", parent:"e0", creation_rank:1,
          content_commit:"h1", payload_handle:"p1" }
}
BINDING {
  subject:"e1", nu_b:"corrections-v4",
  obligations:["amount_matches","second_approval","origin_reference"]
}
CANDIDATE {
  output_commit:"x0", binding_commit:"b0",
  history_head:"H17", nu_k:"adjustment-kernel-v2"
}
VERIFICATION {
  candidate_commit:"x0", binding_commit:"b0", verdict:"FAIL",
  discharge:[["amount_matches","SAT","REQUIRED"],
             ["second_approval","UNSAT","REQUIRED"],
             ["origin_reference","SAT","REQUIRED"]],
  r_block:["second_approval"], r_waived:[]
}
\end{verbatim}
\end{scriptsize}

A second descendant \(e_2\) supplies approval, is rebound and transformed
again, and receives PASS with an empty blocking residual. The new candidate
may share the same output value \(y\) while remaining a distinct candidate
record because its binding and history commitments differ.

\begin{scriptsize}
\begin{verbatim}
COMMIT_CORE {
  id_c:"c0", prior_head:"H17", candidate_commit:"x1",
  verification_commit:"v1", sigma_c:"corrections-v4",
  nonce:"n0", committed_at:"...", authority:"manager:osei",
  r_waived:[]
}
SIGNED_COMMIT { object:"c0", signer:"manager:osei",
                algorithm:"ed25519", signature:"..." }
\end{verbatim}
\end{scriptsize}

Every example record is accompanied in the machine-readable distribution by
its canonical bytes, object identifier, signing preimage, signature, and
expected successor head. The prose form above is explanatory and not itself
the canonical serialization.

\chapter{Failure and Recovery Matrix}
\section{Before external actuation}
Before a transition record is durable, recovery treats the transition as not
having become visible. It may retry the operation as a new attempt, but it must
not claim that the earlier decision occurred. After a complete record is
durable but before canonical application, recovery applies that record
idempotently. It does not rerun POP authentication, REFUSE policy judgment, or
COLLAPSE arbitration against the current world.

For BIND or TRANSFORM, recomputation is permitted only when the corresponding
determinism and historical-input hypotheses hold. VERIFY is replayed from its
signed durable result unless a separate verifier replay theorem applies. A
partial or unauthenticated record is rejected or completed from its
write-ahead transaction; it is never guessed into a valid transition.

\section{Collapse publication}
If a collapse transaction is absent, form a fresh proposal against the current
head. If the transaction is durable and neither \(H\) nor \(L\) reflects it,
apply both. If exactly one side reflects it, apply only the missing side. If
both do, mark the transaction complete. Every application is keyed by the
transaction identifier, preventing duplicate commits or witnesses.

\section{External-effect states}
Before PREPARE is durable, no external request may have been sent. A durable
intent with no attempt witness is pending first transmission. A durable
pre-transmission attempt witness with no conclusive receipt is ambiguous:
recovery queries the external system when possible and retries only under the
declared duplicate-handling contract. A conclusive receipt not yet finalized
causes deterministic FINALIZE and commit publication. Conflicting or duplicate
receipts enter reconciliation; a shared key proves common intent, not that the
external system produced only one effect.

The only stages at which an irreversible effect may already have occurred are
ACTUATE and later. Undeclared side effects during TRANSFORM or VERIFY are
violations of confinement and require their own domain-specific reconciliation.

\chapter{Security Claims and Non-Claims}
\section{Adversarial components}
Against a Byzantine source, the architecture guarantees only that declared
gates and obligations are applied; it does not infer malicious intent from a
policy-compliant request. Kernel and verifier identities make artifact claims
checkable, but do not prove deployment loaded those artifacts without
attestation. History uniqueness holds only under the linearizable arbiter
premise. Clock-relative freshness holds only to the degree the trusted clock
and registries are honest.

Against storage alteration, an authentic later anchor makes rewriting
detectable under the hash assumptions. This is detection, not availability,
truth, or non-equivocation. Against ordinary crashes, checkpoint and record
replay restore canonical internal state under Chapter 23's six hypotheses.
External exactly-once execution remains a non-claim.

\section{General non-claims}
The architecture does not guarantee truth, benevolence, complete evidence,
honest policy, ethical retention, infallible implementation, global ledger
non-equivocation, or adequacy of a physical or institutional model. PASS is
relative warrant. A redaction authorization proves authorization occurred, not
that it was justified. Multiple verifiers provide no multiplicative assurance
against a shared fault without an independence argument. A compensating act
does not reverse the original effect.

\chapter{Research Programme}
\section{Formal and distributed extensions}
Open work includes a verifier-specific replay theorem; mechanized refinement
from the reference service to Appendix B; positive quantitative models of
verifier diversity; causal partial-order history safety; cross-shard atomic
commitment; cross-ledger commitments; and compositional obligation systems.
Ledger non-equivocation requires an extension of the authoritative-head
mechanism from \(H\) to \(L\).

\section{Privacy, policy, and institutions}
Privacy-preserving refusal proofs, completeness proofs for selective ledger
projections, formally governed policy evolution, metadata-minimizing retention,
and empirical methods for detecting institutional capture remain open.
Terminal refusal is represented here as an orthogonal terminal flag, but a
richer algebra of appeal, expiry, jurisdiction, and exceptional reopening may
be preferable.

\section{Integration hypotheses}
An RSVP--Spherepop bridge still needs explicit continuous-to-event observation
maps and consistency laws. KES's per-attempt contraction proposal requires an
independent formalization. Distinction holonomy needs transport operators,
closed-loop structure, and invariants beyond a residue distance. MEM|8 needs a
technical retrieval model proving that admissibility constrains indexing and
scoring before secret material can influence observable rankings.

These are not all “extensions rather than gaps.” Some, such as ledger
non-equivocation, verifier replay, and full implementation refinement, are
unclosed assurance obligations for stronger deployments. The specified core
remains coherent within its stated scope precisely because those stronger
claims are withheld.

\backmatter

\begin{thebibliography}{99}
\addcontentsline{toc}{chapter}{Bibliography}

\subsection*{Distributed Systems and Consensus}

\bibitem{gilbert-lynch}
Seth Gilbert and Nancy Lynch.
``Brewer's Conjecture and the Feasibility of Consistent, Available,
Partition-Tolerant Web Services.''
\emph{ACM SIGACT News} 33, no. 2 (2002): 51--59.
Cited in Chapter 21 for the CAP theorem, with the explicit caution that this
book's safety-versus-liveness tradeoff resembles that result in shape without
having been shown to reduce to it.

\bibitem{lamport-byzantine}
Leslie Lamport, Robert Shostak, and Marshall Pease.
``The Byzantine Generals Problem.''
\emph{ACM Transactions on Programming Languages and Systems} 4, no. 3
(1982): 382--401.
Cited in Chapter 22 as the source of the Byzantine adversary model adopted by
the threat analysis.

\bibitem{lamport-paxos}
Leslie Lamport.
``The Part-Time Parliament.''
\emph{ACM Transactions on Computer Systems} 16, no. 2 (1998): 133--169.
Cited in Chapter 16 as the original Paxos consensus protocol and as a
representative mechanism capable of realizing hypothesis (H2).

\bibitem{ongaro-raft}
Diego Ongaro and John Ousterhout.
``In Search of an Understandable Consensus Algorithm.''
In \emph{Proceedings of the 2014 USENIX Annual Technical Conference},
305--319, 2014.
Cited alongside Paxos in Chapter 16 as an alternative infrastructure capable
of satisfying (H2).

\subsection*{Cryptography}

\bibitem{stevens-sha1}
Marc Stevens, Elie Bursztein, Pierre Karpman, Ange Albertini, and Yarik Markov.
``The First Collision for Full SHA-1.''
In \emph{Advances in Cryptology: CRYPTO 2017}, Lecture Notes in Computer
Science 10401, 570--596, 2017.
The SHAttered result is cited in Chapter 29 as historical confirmation that
collision resistance is an imported and revisable assumption rather than a
merely theoretical caveat.

\bibitem{katz-lindell}
Jonathan Katz and Yehuda Lindell.
\emph{Introduction to Modern Cryptography}. 3rd ed.
Boca Raton, FL: Chapman and Hall/CRC, 2020.
Cited where the book invokes standard cryptographic assumptions, including
collision resistance and existential unforgeability under chosen-message
attack, without reproving them.

\subsection*{This Author's Related Work}

The following unpublished manuscripts are cited in Part VII as sources of its
integration hypotheses. None is re-derived or independently validated here.
Each citation marks the boundary between what this book establishes and what
it borrows from elsewhere in the author's corpus.

\bibitem{flyxion-holonomy}
Flyxion.
``Distinction Holonomy.'' Unpublished manuscript.
Source of the term and general framing adopted, at a deliberately reduced
level of formal commitment, in Chapter 31.

\bibitem{flyxion-wave-machine}
Flyxion.
``Against the Wave Machine.'' Unpublished manuscript.
Source of the MEM|8 resonant-recall proposal discussed as a candidate model,
rather than a validated mechanism, in Chapter 32.

\bibitem{flyxion-conscious-gaia}
Flyxion.
``MEM|8 and the Possibility of a Conscious Gaia.'' Unpublished manuscript.
A more speculative extension of MEM|8 that is not relied upon by this book and
is listed so that readers do not conflate it with the immediate source of
Chapter 32's recall model.

\bibitem{flyxion-rsvp-kes}
Flyxion.
Unpublished manuscripts on the Relativistic Scalar-Vector Plenum (RSVP) and
Kinetic-Event Synthesis (KES).
Sources of the bridge principle discussed in Chapter 30, including the
standing caution that KES's detailed mathematics remains a provisional
reconstruction rather than settled content.

\bibitem{flyxion-spherepop}
Flyxion.
\emph{Spherepop Specification and Reference Implementation Materials}.
Multiple unpublished manuscripts and a Rust reference kernel.
These materials constitute the broader research programme from which the
book's six-transition pipeline is drawn. This book is a focused elaboration of
one part of that programme, not a replacement for it.

\subsection*{A Note on Citation Practice in This Book}

Most of the book's argumentative weight rests on results proved within its own
chapters and cross-referenced internally. External citations appear where the
book imports a genuinely external result, a cryptographic assumption, a named
consensus protocol, a documented historical event, or separately developed
work whose claims are not rechecked here. A reader examining a particular
claim should therefore expect an internal chapter reference in most cases and
an entry from this bibliography only at the explicitly cited boundaries.

\end{thebibliography}

\chapter*{Drafting Sequence}
\addcontentsline{toc}{chapter}{Drafting Sequence}

The recommended drafting order follows dependency rather than final reading
order. First write Chapters 1--4 to stabilize the distinctions and vocabulary.
Next write Chapters 5--9 to define the event types and obligation system.
Then write Chapters 10--19 to complete execution, verification, collapse, and
the witness plane. Draft the reference data model and operational-semantics
appendix immediately afterward so that informal drift becomes visible. Only
then write the implementation, evaluation, and integration parts. The ethical
and institutional chapters should be written last, once the technical limits
of the analogy are explicit.

\chapter*{Candidate Alternative Titles}
\addcontentsline{toc}{chapter}{Candidate Alternative Titles}

\begin{description}[style=nextline]
  \item[Witness Before Commitment]
  Event Memory and the Architecture of Warrant
  \item[Execution Is Not Admission]
  A Calculus of Refusal, Verification, and Irreversible Action
  \item[The Admissible Event]
  History, Constraint, and Commitment in Spherepop
  \item[Memory Without Permission]
  Refusal Ledgers and the Boundary Between Evidence and Action
  \item[Collapse Under Warrant]
  The Spherepop Architecture of Accountable Commitment
\end{description}

\end{document}
